%%%%%%%%%%%%%%%%%%%%%%%%%%%%%%%%%%%%%%%%%%%%%%%%%%%%%%%%%%%%%%%%%
% Contents: The transactions chapter
% $Id: grisbi-manuel-transactions.tex,v 0.4 2002/10/27 Daniel Cartron
% $Id: grisbi-manuel-transactions.tex, v 0.5.0 2004/06/01 Loic Breilloux
% some of its content was in tips chapter: 
% $Id: grisbi-manuel-tips.tex, v 0.4 2002/10/27 Daniel Cartron
% $Id: grisbi-manuel-transactions.tex, v 0.6.0 2011/11/17 Jean-Luc Duflot
% some of its content was in tips chapter:
% $Id: grisbi-manuel-tips.tex, v 0.5.0 2004/06/01 Loic Breilloux
% $Id: grisbi-manuel-transactions.tex, v 0.8.9 2012/04/27 Jean-Luc Duflot
% $Id: grisbi-manuel-transactions.tex, v 1.0 2014/02/12 Jean-Luc Duflot
% some of its content was in accounts chapter:
% $Id: grisbi-manuel-accounts.tex, v 0.8.9 2012/04/27 Jean-Luc Duflot
% $Id: 08-grisbi-manuel-transaction-fr.tex, v 3.0 2026/05 Dominique Brochard
%%%%%%%%%%%%%%%%%%%%%%%%%%%%%%%%%%%%%%%%%%%%%%%%%%%%%%%%%%%%%%%%%

\chapter{Account transactions\label{transactions}}
%\chapter{Opérations d'un compte\label{transactions}}

The various accounts you have created contain transactions. To view the list of accounts stored in your accounts file, use the information bar or the navigation panel (see the \vref{home}, \menus{Home} chapter and the \vref{accounts-list}, \menus{Account List} section).
%Les différents comptes que vous avez créés contiennent les opérations. Pour afficher la liste des comptes enregistrés dans votre fichier de comptes, utilisez la barre d'information ou le panneau de navigation (voir le chapitre \vref{home}, \menus{Accueil} et la section \vref{accounts-list}, \menus{Liste des comptes}).

To view the transactions for an account, select that account from the list: the \menus{Transactions} tab\refimage{transactions-tab-img} will appear by default. It contains four main sections:
%Pour afficher les opérations d'un compte, sélectionnez ce compte dans la liste: l'onglet \menus{Opérations} s'affiche par défaut. Il possède quatre éléments principaux:

\begin{enumerate}
	\item the toolbar (section \vref{transactions-functions});
	%\item la barre d'outils;
	\item the list of transactions (section \vref{transactions-list});
	%\item la liste des opérations;
	\item the transaction entry form (section \vref{transactions-form});
	%\item le formulaire de saisie des opérations;
	\item the \indexword{reconciled balance}\index{balance!reconciled} (section \vref{transactions-balance}).
	%\item le \indexword{solde rapproché}\index{solde!rapproché}.
\end{enumerate}

% espace avant Attention ou Note: 5 mm
\vspacepdf{3mm}
\Note{}: if you have never carried out a reconciliation – i.e. you have not performed any reconciliation entries – Grisbi cannot calculate or display a reconciled balance.
%\Note{}: si vous n'avez jamais fait de rapprochement, donc pas de pointages, Grisbi ne peut ni calculer ni afficher de solde rapproché.

\vspacepdf{3mm}
% image centrée 
\begin{figure}[ht]
	\begin{center}
		\includegraphics[width=0.95\textwidth]{image/screenshot/transactions_tab}
	\end{center}
	\caption{The \menus{Transactions} tab of an account}
	%\caption{Liste des opérations}
	\label{transactions-tab-img}
\end{figure}
% image centrée 


\section{Toolbar\label{transactions-functions}}
%\section{Barre d'outils\label{transactions-functions}}

The toolbar is located at the top of the details panel and offers the following functions:
%La barre d'outils est présente en haut du panneau des détails et présente les fonctions suivantes:

\begin{itemize}
	\item \menus{New transaction}: goes to the end of the list of transactions and opens a blank entry form;
	%\item \menus{Nouvelle Opération}: va à la fin de la liste des opérations et ouvre le formulaire de saisie vierge;
	\item \menus{Delete}: deletes the selected transaction;
	%\item \menus{Supprimer}: supprime l'opération sélectionnée;
	\item \menus{Edit}: opens the input form for the selected transaction and allows you to edit it;
	%\item \menus{Éditer}: ouvre le formulaire de saisie pour l'opération sélectionnée et permet de la modifier;
	\item \menus{Reconcile}: opens the reconciliation page for the selected account (see the section \vref{reconciliation}, \menus{Bank reconciliation});
	%\item \menus{Rapprocher}: ouvre la page de rapprochement du compte sélectionné (voir le chapitre \vref{reconciliation}, \menus{Rapprochement bancaire});
	\item \menus{Print}: opens the window for selecting the printer and its options;
	%\item \menus{Imprimer}: ouvre la fenêtre de sélection de l'imprimante et de ses options;
	\item \menus{View}: opens a drop-down menu; you can choose the \indexword{view of transactions}\index{view!transactions} from four options:
	%\item \menus{Affichage}: ouvre un menu déroulant; vous pouvez choisir l'\indexword{affichage des opérations}\index{affichage!opérations} parmi quatre possibilités:
		\begin{itemize} 
	 		\item \menus{Simple view},
	 		%\item \menus{Vue simple},
	 		\item \menus{Two lines view},
	 		%\item \menus{Mode \frquote{deux lignes}},
	 		\item \menus{Three lines view},
	 		%\item \menus{Mode \frquote{trois lignes}},
	 		\item \menus{Complete view},
	 		%\item \menus{Vue complète},
		 	\item[] \hspace{-15pt} as well as:
		 	%\item[] \hspace{-15pt} ainsi que:
	 		\item whether or not to show \indexword{reconciled operations}\index{show!reconciled operations},
	 		%\item l'affichage ou non des \indexword{opérations rapprochées}\index{affichage!opérations rapprochées},
	 		\item whether or not to show the \indexword{archives lines}\index{display!archives lines} (\vref{datamanagement-archive-display}).
	 		%\item l'affichage ou non des \indexword{lignes d'archives}\index{affichage!lignes d'archives} (\vref{datamanagement-archive-display}).
		\end{itemize}
		\vspacepdf{2mm}
		\Note{}: the options in the \menus{View} menu affect the display of all accounts by default. You may prefer a different display for each account; this can be configured in the \menus{Edit - Preferences - Transactions - List Behaviour} menu (see the section \vref{setup-operations-list-differentiation}, \menus{Account Differentiation}).
		%\Note{}: les options du menu \menus{Affichage} influent par défaut sur l'affichage de tous les comptes. Vous pouvez préférer un affichage différent pour chaque compte, cela est configurable dans le menu \menus{Éditer - Préférences - Opérations - Comportement de la liste} (voir le paragraphe \vref{setup-operations-list-differenciation}, \menus{Différenciation des comptes}).	
	\item a \menus{Search} box, which must contain at least three characters,
	%\item un champ \menus{Rechercher}, à remplir avec trois caractères minimum, 
	\item a drop-down list on the right that allows you to filter the displayed operations by:
	%\item une liste déroulante associée, à droite, permettant de filtrer les opérations affichées selon:
		\begin{itemize} 
			\item the \menus{Payee name},
			%\item le \menus{Nom du tiers},
			\item the \menus{Note},
			%\item la \menus{Remarque},
			\item the \menus{Payee and note} (selected by default),
			%\item les \menus{Tiers et remarque} (sélectionné par défaut),
			\item the \menus{Budgetary line},
			%\item l'\menus{Imputation budgétaire},
			\item the \menus{Category}.
			%\item la \menus{Catégorie}.
		\end{itemize}
		\vspacepdf{2mm}
		\Note{}: the search is case-insensitive and must include at least three (3) alphanumeric characters; it displays only the matching operations in the list of operations. The search is performed ONLY on the displayed operations (note whether or not the \indexword{reconciled operations}\index{show!reconciled operations} are shown) and does not include the archives.
		%\Note{}: la recherche est insensible à la casse et doit se faire avec un minimum de trois (3) caractères alpha-numériques; elle affiche uniquement toutes les opérations correspondantes dans la liste des opérations. La recherche ne se fait QUE sur les opérations affichées (attention à l'affichage ou non des \indexword{opérations rapprochées}\index{affichage!opérations rapprochées}) et ne concerne pas les archives.
\end{itemize}


\section{List of transactions\label{transactions-list}}
%\section{Liste des opérations\label{transactions-list}}


\subsection{Description\label{transactions-list-description}}
%\subsection{Description\label{transactions-list-description}}

The list of transactions is displayed in the details panel\refimage{transactions-tab-img} below the toolbar.
%La liste des opérations s'affiche dans le panneau des détails\refimage{transactions-list-img} sous la barre d'outils.

The column header bar is displayed at the top. You can \indexword{widen or narrow a column}\index{column!width} by clicking on the separator between two columns and dragging it. To restore the columns to their default width, select \menus{View - Reset the column width} from the menu bar.
%Elle affiche en haut la barre des libellés des colonnes. Vous pouvez \indexword{élargir ou rétrécir une colonne}\index{colonne!largeur} en cliquant sur le séparateur entre deux colonnes et en le déplaçant. Pour rétablir la largeur des colonnes à leur valeur par défaut, sélectionnez dans la barre de menu \menus{Affichage - Réinitialiser la largeur des colonnes}.

You can scroll the list of transactions up or down using the mouse wheel, or by dragging the vertical scroll bar. To scroll left or right, use the mouse to drag the horizontal scroll bar.
%Vous pouvez déplacer la liste des opérations vers le haut ou vers le bas avec la molette de la souris, ou bien avec la souris et l'ascenseur vertical. Le déplacement vers la gauche ou la droite se fait avec la souris et l'ascenseur horizontal.

You can choose whether or not to display:
%Vous pouvez choisir d'afficher ou non:
\begin{itemize}
	\item the \indexword{reconciled transactions}\index{display!reconciled transactions} by selecting:
	%\item les \indexword{opérations rapprochées}\index{affichage!opérations rapprochées} en sélectionnant:
		\begin{itemize}
			\item in the menu bar \menus{View - Show reconciled transactions} (\keys{Alt+R}) (see the section \vref{home-menus-display}),
			%\item dans la barre de menu \menus{Affichage - Montrer les opérations rapprochées} (\keys{Alt+R}) (voir la section \vref{home-menus-display}),
			\item in the toolbar \menus{View - Show reconciled Transactions} (section \vref{transactions-functions});
			%\item dans la barre d'outils \menus{Affichage - Montrer les opérations rapprochées} (section \vref{transactions-functions});
		\end{itemize} 
	\item the \indexword{lines archives}\index{display!lines archives} (see the section \menus{Archives display} \vref{datamanagement-archive-display}):
	%\item les \indexword{lignes d'archives}\index{affichage!lignes d'archives} (voir la section \menus{Affichage des archives} \vref{datamanagement-archive-display}):
		\begin{itemize}
			\item in the menu bar \menus{View - Show lines archives} (\keys{Alt+L}) (see the section \vref{home-menus-display}),
			%\item dans la barre de menu \menus{Affichage - Montrer les lignes d'archives} (\keys{Alt+L}) (voir la section \vref{home-menus-display}),
			\item in the toolbar \menus{View - Show lines archives} (section \vref{transactions-functions}).
			%\item dans la barre d'outils \menus{Affichage - Montrer les lignes d'archives} (section \vref{transactions-functions}).
		\end{itemize}
\end{itemize}


% espace pour changement de thème
\vspacepdf{3mm}
From the perspective of the Grisbi database, a transaction consists of its \indexword{information fields}\index{information field!definition}. The information fields for each transaction are displayed in the various \indexword{cells}\index{cell!information field} defined by the intersection of rows and columns. There can be a maximum of four rows and seven columns, which defines a maximum of twenty-eight information fields for each transaction.
%Du point de vue de la base de données de Grisbi, une opération est constituée de ses \indexword{champs d'information}\index{champ d'information!définition}. Les champs d'information de chaque opération s'affichent dans les différentes \indexword{cellules}\index{cellule!champ d'information} définies par l'intersection des lignes et des colonnes. Il ne peut y avoir au maximum que quatre lignes et sept colonnes, ce qui définit un maximum de vingt-huit champs d'information pour chaque opération.

Each transaction can be displayed over 1, 2, 3 or 4 lines, depending on the selected display mode (see the \vref{home-menus-display} section, the \menus{View} menu in the menu bar, or the \menus{View} function in the \vref{transactions-functions}, \menus{Toolbar} section). The current display is indicated by a tick.
%Chaque opération peut être affichée sur 1, 2, 3 ou 4 lignes, suivant le mode d'affichage sélectionné (voir la section \vref{home-menus-display}, menu \menus{Affichage} dans la barre de menu ou la fonction \menus{Affichage} dans la section \vref{transactions-functions}, \menus{Barre d'outils}). L'affichage courant y est indiqué par une coche. 

You can also configure the content of these lines, as well as the saving of display settings, individually for each account, in the \menus{Edit - Preferences - Transactions} (see the sections \vref{setup-operations-list}, \menus{List behavior} and \vref{setup-operations-cells}, \menus{Transactions list cells}).
%Vous pouvez aussi configurer le contenu de ces lignes, ainsi que la mémorisation des réglages de l'affichage, individuellement pour chaque compte, dans le menu \menus{Éditer - Préférences - Opérations} (voir les sections \vref{setup-operations-list}, \menus{Comportement de la liste} et \vref{setup-operations-cells}, \menus{Cellules de la liste des opérations}).

The \indexword{column label}\index{column!label} is always the name of the information field displayed in the first row of the transactions.
%Le \indexword{libellé d'une colonne}\index{colonne!libellé} est toujours le nom du champ d'information affiché dans la première ligne des opérations.

To improve the readability of the display, Grisbi does not show any borders around cells, columns, rows or transactions, but instead uses alternating colours of  \fboxsep=2pt{\colorbox[RGB]{215,215,255}{purple background}} and white for each line.
%Pour des raisons de lisibilité de l'affichage, Grisbi n'affiche aucune bordure aux cellules, colonnes, lignes et opérations, mais présente une alternance de couleurs de  \fboxsep=3pt{\colorbox[RGB]{215,215,255}{fond violet}} et blanc à chaque ligne.


\subsection{Information and input fields\label{transactions-list-fields}}
%\subsection{Champs d'information et de saisie\label{transactions-list-fields}}

Grisbi can manage, for each transaction, a complete list of information fields and input fields\index{information fields}\index{input fields}, which are as follows:
%Grisbi sait gérer, pour chaque opération, une liste complète de champs d'information et de saisie\index{champs d'information}\index{champs de saisie}, qui sont les suivants:

\begin{itemize}
	\item \menus{Date} of the transaction\index{date!transaction}: \indexword{date} on which the transaction took place; this can be used to calculate balances;
	%\item \menus{Date} de l'opération\index{date!opération}: \indexword{date} à laquelle est réalisée l'opération; elle peut être utilisée pour calculer les soldes;
	\item \menus{Value date}\index{date!value}: \indexword{date} on which the bank records the transaction; this is the date used for calculating interest and fees; it may be used to calculate balances;
	%\item \menus{Date de valeur}\index{date!valeur}: \indexword{date} à laquelle la banque comptabilise l'opération; c'est cette date qui fait foi pour le calcul des intérêts et des agios; elle peut être utilisée pour calculer les soldes;
	\item \menus{Payee};
	%\item \menus{Tiers};
	\item \menus{Budgetary line};
	%\item \menus{Imputation budgétaire};
	\item \menus{Debit}: amount debited from the account balance; in the case of transactions in a \indexword{different currency}\index{currency!different} from that of the account, the transaction amount in the different currency;
	%\item \menus{Débit}: montant débité du solde du compte; dans le cas d'opérations dans une \indexword{devise différente}\index{devise!différente} de celle du compte, montant de l'opération dans la devise différente;
	\item \menus{Credit};
	%\item \menus{Crédit};
	\item \menus{Balance}: an amount always calculated \emph{after} the operation displayed on the same line; it depends chronologically on the selected \indexword{primary sort}\index{sort!primary}; a \indexword{negative balance}\index{balance!negative} is displayed in \textcolor{red}{red};
	%\item \menus{Solde}: montant toujours calculé \emph{après} l'opération affichée sur la même ligne; il dépend chronologiquement du \indexword{tri primaire}\index{tri!primaire} sélectionné; un \indexword{solde négatif}\index{solde!négatif} est affiché en \textcolor{red}{rouge};
	\item \menus{Amount}: for transactions in a \indexword{different currency}\index{currency!different} from that of the account, the transaction amount in the account’s currency;
	%\item \menus{Montant}: dans le cas d'opérations dans une \indexword{devise différente}\index{devise!différente} de celle du compte, montant de l'opération dans la devise du compte;
	\item \menus{Method of payment};
	%\item \menus{Moyen de paiement};
	\item \menus{Reconciliation reference}: the number of each reconciliation entry in each account;
	%\item \menus{\No rapprochement}: numéro de chaque rapprochement dans chaque compte;
	\item \menus{Financial year};
	%\item \menus{Exercice};
	\item \menus{Category};
	%\item \menus{Catégorie};
	\item \menus{P/R}: transaction status, \emph{P (Pointed)} or \emph{R (Reconciled)};
	%\item \menus{P/R}: statut de l'opération, \emph{P (Pointé)} ou \emph{R (Rapproché)};
	\item \menus{Voucher}, for example the invoice number (\menus{Voucher} and \menus{Transaction number} are cross-references that ensure the traceability of the transaction);
	%\item \menus{Pièce comptable}, par exemple le numéro de la facture (\menus{Pièce comptable} et \menus{Numéro d'opération} sont des références réciproques qui assurent la traçabilité de l'opération);
	\item \menus{Notes}: any relevant information about the transaction;
	%\item \menus{Remarques}: informations éventuelles sur l'opération;
	\item \menus{Bank references}, only relevant if you are recording a cheque you have received;
	%\item \menus{Infos banque/guichet}, utile uniquement si vous enregistrez un chèque reçu;
	\item \menus{Transaction number}: this is a unique serial number that Grisbi assigns to each transaction; you can transfer this number to the corresponding voucher (\menus{Transaction number} and \menus{Voucher} are cross-references that ensure the transaction can be traced);
	%\item \menus{\No opération}: il s'agit d'un numéro d'ordre unique que Grisbi attribue à chaque opération; vous pouvez reporter ce numéro sur la pièce comptable correspondante (\menus{\No opération} et \menus{Pièce comptable} sont des références réciproques qui assurent la traçabilité de l'opération);
	\item \menus{Cheque number} or \menus{Cheque/Transfer number}, if the selected payment method allows this information to be entered;
	%\item \menus{\No de chèque} ou \menus{\No chèque/virement}, si le mode de règlement sélectionné permet cette saisie;
	\item \menus{Currency} of the transaction;
	%\item \menus{Devise} de l'opération;
	\item \menus{Change} corresponding to the exchange rate;
	%\item \menus{Change} correspondant au taux de change;
	\item \menus{\indexword{Contra-transaction}\index{Contra-transaction} method of payment}: you can configure this for each account in the \menus{Edit - Preferences - Resources} menu (see the section \vref{setup-resources-modes}, \menus{Payment methods});
	%\item \menus{Mode de règlement de la \indexword{contre-opération}}\index{contre-opération}: vous pouvez le configurer pour chaque compte dans le menu \menus{Éditer - Préférences - Ressources} (voir la section \vref{setup-resources-modes}, \menus{Modes de règlement});
	\item \menus{Automatic/Manual}: status of a scheduled transaction's execution.
	%\item \menus{Auto/Manuel}: statut de l'exécution d'une opération planifiée
\end{itemize}

% espace pour changement de thème
\vspacepdf{3mm}
The list of transactions and the transaction/scheduled form each use a subset of these fields. You can view the two complete lists in the tables under the \menus{Edit - Preferences} menu:
%La liste des opérations et le formulaire de saisie utilisent chacun un sous-ensemble de ces champs. Vous pouvez voir les deux listes complètes dans les tableaux du menu \menus{Éditer - Préférences}:
\begin{itemize}
	\item section \menus{Transactions - Transactions list cells} (see section \vref{setup-operations-cells} \menus{Cells in the transaction list});
	%\item section \menus{Opérations - Cellules de la liste des opérations} (voir la section \vref{setup-operations-cells} \menus{Cellules de la liste des opérations});
	\item section \menus{Transaction Form - Content} (see section \vref{setup-form-content} \menus{Content}).
	%\item section \menus{Formulaire des opérations - Contenu} (voir la section \vref{setup-form-content} \menus{Contenu}).
\end{itemize}


\subsection{Management of information fields\label{transactions-list-fields-manage}}
%\subsection{Gestion des champs d'information\label{transactions-list-fields-manage}}

By default, Grisbi displays a number of these fields (payee, category, debit, etc.) in the transaction list. However, you can also \indexword{manage the information fields}\index{information field!management} differently, i.e. display new fields, modify them or hide them, as you see fit.
%Grisbi affiche par défaut, dans la liste des opérations, un certain nombre de ces champs (tiers, catégorie, débit, etc.). Mais vous pouvez aussi \indexword{gérer les champs d'information}\index{champ d'information!gestion} différemment, c'est à dire afficher de nouveaux champs, les modifier ou en supprimer l'affichage, à votre convenance.
You can either carry out these actions in the \menus{Edit - Preferences - Transactions - Transactions list cells}, which provides a complete overview of the layout of fields within the cells and allows you to make several changes easily (see the section \vref{setup-operations-cells}, \menus{Transactions list cells}), or do so directly in the transaction list, using the methods below:
%Vous pouvez soit faire ces manipulations dans le menu \menus{Éditer - Préférences - Opérations - Cellules de la liste des opérations}, qui offre un aperçu complet de la disposition des champs dans les cellules, et permet de faire plusieurs modifications facilement (voir la section \vref{setup-operations-cells}, \menus{Cellules de la liste des opérations}), soit le faire directement dans la liste des opérations, en suivant les méthodes ci-dessous:


\subsubsection{Adding an information field\label{transactions-list-fields-add}}
%\subsubsection{Ajout d'un champ d'information\label{transactions-list-fields-add}}

To \indexword{add an information field}\index{information field!add} to a cell, do the following:
%Pour \indexword{ajouter un champ d'information}\index{champ d'information!ajouter} dans une cellule, procédez comme suit:

\begin{enumerate}
	\item select \menus{View - Complete view} from the toolbar, or \menus{View - Show four lines per transaction} from the menu bar, to ensure that you can see all twenty-eight cells;
	%\item sélectionnez dans la barre d'outils \menus{Affichage - Vue complète} ou dans la barre de menu \menus{Affichage - Montrer quatre lignes par opération}, pour vous assurer de voir toutes les vingt-huit cellules;
	\item \texttt{right-click} in one of the empty cells in the column where you want to add this field: be precise, as there is no way of telling which cell you have clicked on;
	%\item cliquez-droit dans une des cellules vides de la colonne où vous voulez ajouter ce champ: faites-le précisément, car rien ne permet d'identifier sur quelle cellule vous avez cliqué;
	\item a context menu will appear; select \menus{Change cell content}; check that there is no information field already in this cell by ensuring that none of the fields in the drop-down list has a tick to the left of its name;
	%\item un menu contextuel s'affiche; sélectionnez \menus{Modifier le contenu de la cellule}; vérifiez qu'il n'y a pas déjà de champ d'information dans cette cellule, en constatant qu'aucun champ de la liste déroulante n'a de coche à gauche de son nom;
	\item select an information field from the list; this field will be added to the selected cell, and all transactions in the list will display it.
	%\item choisissez un champ d'information proposé dans la liste; ce champ s'ajoute dans la cellule choisie, et toutes les opérations de la liste l'affichent.
\end{enumerate}


\subsubsection{Editing an information field\label{transactions-list-fields-modify}}
%\subsubsection{Modification d'un champ d'information\label{transactions-list-fields-modify}}

To \indexword{edit an information field}\index{information field!edit} in a cell, do the following:
%Pour \indexword{modifier un champ d'information}\index{champ d'information!modifier} dans une cellule, procédez comme suit:

\begin{enumerate}
	\item select \menus{View - Complete view} from the toolbar, or \menus{View - Show four lines per transaction} from the menu bar, to ensure that you can see all twenty-eight cells;
	%\item sélectionnez \menus{Affichage - Vue complète} dans la barre d'outils ou \menus{Affichage - Montrer quatre lignes par opération} dans la barre de menu, pour vous assurer de voir toutes les vingt-huit cellules;
	\item \texttt{right-click} in one of the empty cells in the column where you want to modify this field: be precise, as there is no way of telling which cell you have clicked on;
	%\item cliquez-droit dans une des cellules de la colonne dont vous voulez modifier le champ: faites-le précisément, car rien ne permet d'identifier sur quelle cellule vous avez cliqué;
	\item a context menu will appear; select \menus{Change cell content}; the current information field in this cell is indicated by a tick to the left of its name;
	%\item un menu contextuel s'affiche; sélectionnez \menus{Modifier le contenu de la cellule}; le champ d'information actuel dans cette cellule est indiqué par une coche à gauche de son nom;
	\item select another field from the list; this field will replace the previous one in the selected cell, and all transactions in the list will display it.
	%\item choisissez un autre champ d'information proposé dans la liste; ce champ remplace le précédent dans la cellule choisie, et toutes les opérations de la liste l'affichent.
\end{enumerate}


\subsubsection{Deleting an information field\label{transactions-list-fields-remove}}
%\subsubsection{Suppression d'un champ d'information\label{transactions-list-fields-remove}}

To \indexword{delete an information field}\index{information field!delete} in a cell, do the following:
%Pour \indexword{supprimer un champ d'information}\index{champ d'information!supprimer} dans une cellule, procédez comme suit:

\begin{enumerate}
	\item select \menus{View - Complete view} from the toolbar, or \menus{View - Show four lines per transaction} from the menu bar, to ensure that you can see all twenty-eight cells;
	%\item sélectionnez \menus{Affichage - Vue complète} dans la barre d'outils ou \menus{Affichage - Montrer quatre lignes par opération} dans la barre de menu, pour vous assurer de voir toutes les vingt-huit cellules;
	\item \texttt{right-click} in the cell of the field you want to delete: be sure to do this carefully, as there is no way of telling which cell you have clicked on;
	%\item cliquez-droit dans la cellule du champ à supprimer: faites-le précisément, car rien ne permet d'identifier sur quelle cellule vous avez cliqué;
	\item a context menu will appear; select \menus{Change cell content}; the current information field in this cell is indicated by a tick to the left of its name;
	%\item un menu contextuel s'affiche; sélectionnez \menus{Modifier le contenu de la cellule}; le champ d'information actuel dans cette cellule est indiqué par une coche à gauche de son nom;
	\item select \menus{Clear Cell} from the list; the current information field will disappear for all operations in the list (but the information will not be deleted from your accounts file under any circumstances).
	%\item choisissez \menus{Effacer la cellule} dans la liste; le champ d'information actuel disparaît pour toutes les opérations de la liste (mais les informations ne sont en aucun cas supprimées de votre fichier de comptes).
\end{enumerate}


\subsubsection{Moving an information field\label{transactions-list-fields-move}}
%\subsubsection{Déplacement d'un champ d'information\label{transactions-list-fields-move}}

To \indexword{move an information field}\index{information field!move}, you can delete it and then add it to another empty field, but you can also do this via the \menus{Edit - Preferences - Transactions - Transactions list cells} (see the section \vref{setup-operations-cells}, \menus{Transactions list cells}).
%Pour \indexword{déplacer un champ d'information}\index{champ d'information!déplacer}, vous pouvez le supprimer puis l'ajouter dans un autre champ vide, mais vous pouvez aussi le faire dans le menu \menus{Éditer - Préférences - Opérations - Cellules de la liste des opérations} (voir la section \vref{setup-operations-cells}, \menus{Cellules de la liste des opérations}).


\subsection{Sorts\label{transactions-list-sorts}}
%\subsection{Tris\label{transactions-list-sorts}}

To change the order in which transactions are displayed, you can \indexword{sort the transactions}\index{sort!transactions} by one of the information fields; proceed as follows:
%Pour changer l'ordre d'affichage des opérations, vous pouvez faire un \indexword{tri sur les opérations}\index{tri!opérations} sur un des champs d'information; procédez comme suit:

\begin{enumerate}
	\item \texttt{right-click} on the column header containing the field selected as the \gls{sorting} criterion;
	%\item cliquez-droit sur le libellé de la colonne qui contient le champ d'information choisi comme critère de \gls{tri};
	\item a context menu displays a list of the fields in that column; select the field you wish to use as the sort criterion;
	%\item un menu contextuel affiche la liste des champs de cette colonne; sélectionnez le champ choisi comme critère de tri;
	\item all transactions are displayed, sorted according to the selected criterion;
	%\item toutes les opérations s'affichent, triées selon le critère choisi;
	\item a small black triangle appears to the right or left of the column heading; it points downwards if the sort is ascending, and upwards if it is descending;
	%\item un petit triangle noir apparaît à droite ou à gauche du libellé de la colonne; il pointe vers le bas si le tri est ascendant, vers le haut si le tri est descendant;

% saut de ligne pour indentation correcte de la note dans la liste
	\Note{}: depending on the theme of your desktop environment or the window manager you are using, these triangles can be replaced with other characters such as +, -, >, <, etc.
	%\Note{}: ces triangles peuvent être remplacés, en fonction du thème de l'environnement de bureau ou du gestionnaire de fenêtres que vous utilisez, par d'autres caractères tels que +, -, >, <, etc.
	\item click on this triangle to change the sort order.
	%\item cliquez sur ce triangle pour changer l'ordre du tri.
\end{enumerate}

\Note{}: sorting the display in this way does not affect the balance for each line, which is calculated after each transaction; depending on the sort order and display options you have chosen, the balance for each line may not appear in chronological order.
%\Note{}: un tri sur l'affichage réalisé ainsi n'affecte en aucun cas le solde à chaque ligne, qui est calculé après chaque opération; suivant le tri et les options d'affichage que vous avez choisi, le solde à chaque ligne peut présenter une évolution non chronologique.

%espace pour changement de thème
\vspacepdf{3mm}
You can also configure the general sort order in the \menus{Edit - Preferences} menu (see the section \vref{setup-operations-list-modes}, \menus{Transaction list behavior}).
%Vous pouvez aussi configurer l'ordre général des tris dans le menu \menus{Éditer - Préférences} (voir la section \vref{setup-operations-list-modes}, \menus{Comportement de la liste des opérations}).

\vspacepdf{3mm}
\Note{}: the choice of the \gls{primary sorting} criterion (transaction date or value date) that may have been selected there will necessarily affect the balance displayed on each transaction line, as Grisbi calculates and displays the balance on each line.
%\Note{}: le choix du critère de \gls{tri primaire} (date d'opération ou date de valeur) qui peut y avoir été fait modifie nécessairement le solde affiché à chaque ligne d'opération, puisque Grisbi calcule et affiche le solde à chaque ligne.


\section{Transaction/Scheduled entry form\label{transactions-form}}
%\section{Formulaire de saisie\label{transactions-form}}

The entry form is located below the list of transactions\refimage{transactions-form-img}.
%Le formulaire de saisie est situé en-dessous de la liste des opérations\refimage{transactions-form-img}.

% image centrée 
\begin{figure}[htbp]
	\begin{center}
	\includegraphics[width=0.95\textwidth]{image/screenshot/transactions_form}
	\end{center}
	\caption{Customised transaction entry form}
	%\caption{Formulaire personnalisé de saisie des opérations}
	\label{transactions-form-img}
\end{figure}
% image centrée 


\subsection{Form visibility\label{transactions-form-exposure}}
%\subsection{Visibilité du formulaire\label{transactions-form-exposure}}

The form can be hidden or displayed as required.
%Le formulaire peut être masqué ou affiché selon vos besoins.


\subsubsection{Displaying the form\label{transactions-form-exposure-display}}
%\subsubsection{Affichage du formulaire\label{transactions-form-exposure-display}}

To \indexword{open or display}\index{entry form!display} the entry form, use one of the following methods:
%Pour \indexword{ouvrir ou afficher}\index{formulaire de saisie!afficher} le formulaire de saisie, utilisez une des méthodes suivantes:

\begin{itemize}
	\item in the menu bar, select \menus{Edit - New transaction} if an account is selected;
	%\item dans la barre de menu, sélectionnez \menus{Éditer - Nouvelle opération} si un compte est sélectionné;
	\item in the menu bar, select \menus{View} and tick the box next to the \menus{Show transaction form} option;
	%\item dans la barre de menu, sélectionnez \menus{Affichage} et cochez la case de l'option \menus{Montrer le formulaire de saisie des opérations};
	\item select the blank line at the bottom of the list of transactions and press the \keys{Enter \return} key;
	%\item sélectionnez la ligne vide en bas de la liste des opérations et appuyez sur la touche \keys{Entrée \return};
	\item double-click on the empty line at the bottom of the list of transactions;
	%\item double-cliquez sur la ligne vide en bas de la liste des opérations;
	\item press the \keys{Ctrl+T} key combination;
	%\item appuyez sur la combinaison de touches \keys{Ctrl+T};
	\item in the toolbar, click on \menus{New transaction};
	%\item dans la barre d'outils, cliquez sur \menus{Nouvelle opération};
	\item in the context menu of the transaction list, select \menus{New transaction};
	%\item dans le menu contextuel de la liste des opérations, sélectionnez \menus{Nouvelle opération};
	\item click on the small triangle next to the label \indexword{show/hide} on the transaction entry form to the left of the label.
	%\item cliquez sur le petit triangle d'\indexword{affichage/masquage}\index{formulaire de saisie!masquer-afficher} du formulaire de saisie des opérations à gauche de son libellé.
	
	% saut de ligne pour indentation correcte de la note dans la liste
	\vspacepdf{2mm}
	\Note{}: depending on the theme of your desktop environment or the window manager you are using, these triangles can be replaced with other characters such as +, -, >, <, etc.
	%\Note{}: ce triangle peut être remplacé, en fonction du thème de l'environnement de bureau ou du gestionnaire de fenêtres que vous utilisez, par d'autres caractères tels que +, -, >, <, etc.
\end{itemize}

The last line, which is completely blank, at the bottom of the list of operations, appears in \fboxsep=2pt{\colorbox[RGB]{245,155,155}{dark pink background}}, and the blank input form is displayed, showing only the current date, or the last date entered, in \fboxsep=2pt{\colorbox[RGB]{61,174,233}{\textcolor{white}{\textsf{white on a blue background}}}}.
%La dernière ligne, entièrement vide, en bas de la liste des opérations, apparaît sur \fboxsep=3pt{\colorbox[RGB]{245,155,155}{fond rose foncé}}, et le formulaire de saisie vierge s'affiche, avec seule la date du jour, ou la dernière date entrée, en \fboxsep=3pt{\colorbox[RGB]{61,174,233}{\textcolor{white}{\textsf{blanc sur fond bleu}}}}.

\subsubsection{Hiding the form\label{transactions-new-exposure-hide}}
%\subsection{Masquage du formulaire\label{transactions-new-hide}}

To \indexword{hide}\index{entry form!hide} the entry form, use one of the following methods:
%Pour \indexword{masquer}\index{formulaire de saisie!masquer} le formulaire de saisie, utilisez une des méthodes suivantes:

\begin{itemize}
	\item in the menu bar, select \menus{View} and uncheck the box next to the \menus{Show transaction form} option;
	%\item dans la barre de menu, sélectionnez \menus{Affichage} et décochez la case de l'option \menus{Montrer le formulaire de saisie des opérations};
	\item click on the small triangle next to the label \indexword{show/hide}\index{transaction form!show-hide} on the transaction entry form to the left of the label.
	%\item cliquez sur le petit triangle d'\indexword{affichage/masquage}\index{formulaire de saisie!afficher-masquer} du formulaire de saisie des opérations à gauche de son libellé.
	
	% saut de ligne pour indentation correcte de la note dans la liste
	\vspacepdf{2mm}
	\Note{}: depending on the theme of your desktop environment or the window manager you are using, these triangles can be replaced with other characters such as +, -, >, <, etc.
	%\Note{}: ce triangle peut être remplacé, en fonction du thème de l'environnement de bureau ou du gestionnaire de fenêtres que vous utilisez, par d'autres caractères tels que +, -, >, <, etc.
\end{itemize}


\subsection{Form contents\label{transactions-form-content}}
%\subsection{Contenu du formulaire\label{transactions-form-content}}

The entry form includes the following seven fields by default:
%Le formulaire de saisie comprend par défaut les sept champs suivants:
\begin{itemize}
	\item \menus{Date} of the operation,
	%\item \menus{Date} de l'opération,
	\item \menus{Payee},
	%\item \menus{Tiers},
	\item \menus{Debit},
	%\item \menus{Débit},
	\item \menus{Credit},
	%\item \menus{Crédit},
	\item \menus{Categories : Sub-categories},
	%\item \menus{Catégorie},
	\item \menus{Method of payment}
	%\item \menus{Moyen de paiement},
	\item \menus{Notes}
	%\item \menus{Remarques}.
\end{itemize}
You may need additional fields, such as \menus{Financial year}, \menus{Budgetary line} or others; this is done in the form content settings. There are two ways to access them:
%Vous pouvez avoir besoin de champs supplémentaires, par exemple \menus{Exercice}, \menus{Imputation budgétaire} ou autres; cela se fait dans les paramètres du contenu du formulaire. Vous pouvez y accéder de deux façons différentes:
\begin{itemize}
	\item either by \texttt{right-click} in the entry form, in a grey area outside the entry fields, and then clicking the \menus{Configure the form} button (see the section \vref{setup-form}, \menus{Transaction Form - Content}),
	%\item soit par un clic-droit dans le formulaire de saisie, dans une zone grise en dehors des champs de saisie, puis en cliquant sur le bouton \menus{Configurer le formulaire} (voir la section \vref{setup-form}, \menus{Formulaire des opérations - Contenu}),
	\item or by going to the \menus{Edit - Preferences - Transaction Form - Content} menu to edit the fields as you wish.
	%\item soit en allant dans le menu \menus{Éditer - Préférences - Formulaire des opérations - Contenu} pour modifier les champs à votre convenance.
\end{itemize}
%espace pour changement de thème
\vspacepdf{3mm}
For a full list of the fields that Grisbi can handle, see the section \vref{transactions-list-fields}, \menus{Information and entry fields}
%Pour connaître l'ensemble des champs que peut gérer Grisbi, voir la section \vref{transactions-list-fields}, \menus{Champs d'information et de saisie}.

\vspacepdf{3mm}

Once the form is displayed, a context menu that can be accessed by \texttt{right-click} in an entry field allows you to perform the following actions:
%Une fois le formulaire affiché, un menu contextuel accessible par un clic-droit dans un champ de saisie permet d'effectuer les actions suivantes:
\begin{itemize}
	\item \menus{Cut} (selection);
	%\item \menus{Couper} (la sélection);
	\item \menus{Copy} (selection);
	%\item \menus{Copier} (la sélection);
	\item \menus{Paste} ()selection;
	%\item \menus{Coller} (la sélection);
	\item \menus{Delete} (selection);
	%\item \menus{Supprimer} (la sélection);
	\item[]\hspace{-12pt} then%puis
	\item \menus{Select All} (in the entry field);
	%\item \menus{Tout sélectionner} (dans le champ de saisie);
	\item \menus{Insert Emoji}\index{emoji}.
	%\item \menus{Insérer un émoji}\index{émoji}.
\end{itemize}


\section{Pointed balance\label{transactions-balance}}
%\section{Solde pointé\label{transactions-balance}}

Below the list of transactions, on the same line as the label \menus{Transaction/Scheduled form}, two pieces of information are displayed on the right:
%En-dessous de la liste des opérations, sur la même ligne que le libellé \menus{Formulaire de saisie des opérations}, s'affichent à droite deux informations:
\begin{itemize}
	\item the date of the last reconciled bank statement,
	%\item la date du dernier relevé bancaire rapproché,
	\item the \indexword{reconciled balance}\index{balance!reconciled} corresponding to this latest reconciliation, unless you reconcile your transactions very regularly, in which case it corresponds to the balance as recorded by your bank (see the chapter \vref{reconciliation}, \menus{Bank reconciliation}).
	%\item le \indexword{solde pointé}\index{solde!pointé} correspondant à ce dernier rapprochement, à moins que vous ne pointiez très régulièrement vos opérations, auquel cas il correspond au solde tel que le connaît votre banque (voir le chapitre \vref{reconciliation}, \menus{Rapprochement bancaire}).
\end{itemize}


\section{Selecting an transaction\label{transactions-selection}}
%\section{Sélection d'une opération\label{transactions-selection}}

There are two ways to select an transaction:
%Pour sélectionner une opération, vous avez deux moyens:
\begin{itemize}
	 \item click on one of its lines;
	 %\item cliquez sur une de ses lignes;
	 \item move the selection using the \keys{\arrowkeyup} and \keys{\arrowkeydown} keys.
	 %\item déplacez la sélection avec les touches du clavier \keys{\arrowkeyup} et \keys{\arrowkeydown}.
\end{itemize}

The transaction then appears in \fboxsep=2pt{\colorbox[RGB]{245,155,155}{dark pink background}} and its details are displayed in the form.
%L'opération apparaît alors sur \fboxsep=3pt{\colorbox[RGB]{245,155,155}{fond rose foncé}} et ses détails s'affichent dans le formulaire.

\texttt{Right-click} on a transaction line displays a context menu containing the following functions, depending on the context\refimage{transactions-contextMenu-img}:
%Un clic-droit sur la ligne d'une opération affiche un menu contextuel avec, selon le contexte, les fonctions suivantes:

\vspace{3mm}
\noindent
\begin{minipage}{.65\linewidth}
\begin{itemize}[rightmargin=.6cm]
	\item \menus{Displays the contra-transaction}\index{contra-transaction!display}: this option is only available if the transaction is a transfer;
	%\item \menus{Afficher la contre-opération}\index{contre-opération!afficher}: ce choix est présent seulement si l'opération est un virement;
	\item \menus{Edit transaction};
	%\item \menus{Éditer l'opération};
	\item \menus{New transaction};
	%\item \menus{Nouvelle opération};
	\item \menus{Delete transaction};
	%\item \menus{Supprimer l'opération};
	\item \menus{Use selected transaction as a template};
	%\item \menus{Utiliser l'opération sélectionnée comme modèle};
	\item \menus{Clone transaction};
	%\item \menus{Cloner l'opération};
	\item \menus{Convert transaction to scheduled transaction}: this makes it easy to create due dates retrospectively;
	%\item \menus{Convertir l'opération en opération planifiée}: cela permet de créer facilement des échéances après coup;
	\item \menus{Move transaction to another account};
	%\item \menus{Déplacer l'opération vers un autre compte};
	\item \menus{Search} (see the chapter \vref{search});
	%\item \menus{Rechercher} (voir le chapitre \vref{search});
	\item \menus{Change cell content}.
	%\item \menus{Modifier le contenu de la cellule}.
\end{itemize}
\end{minipage}
\hspace{10pt}
\begin{minipage}{.35\linewidth}
%\begin{wrapfigure}{R}{0.4\textwidth}%50mm
%\vspace{-\intextsep}				% space above the floating (minus intextsep=separation between float and text in text)
	\vspace{-25pt}
	\centering							% centering the floating figure in the "wrap"
%\raggedleft
	\includegraphics[width=1\textwidth]{image/screenshot/transactions_contextMenu}
	\vspace{-15pt}						% space between the floating and the caption below the floating
	\captionsetup{%						% for a change of options of the caption package
	format=plain,					% to avoid error message with labelsep option below, default=hang
	type=figure,%
	name=Fig.,%						% rename label caption, default="Figure" (or "Table")
	justification=centering,		% centring of text and caption label 
	labelsep=newline}				% define the separator between the label and the caption text
	\caption{Context menu for an transaction}		% \caption is mandatory to reference a figure in lof
	%\caption{Menu contextuel d'une opération}		% \caption is mandatory to reference a figure in lof
	\vspace{-10pt}						% space below the caption
	\label{transactions-contextMenu-img}
%\end{wrapfigure}
\end{minipage}


\section{Entering a new transaction\label{transactions-new}}
%\section{Saisie d'une nouvelle opération\label{transactions-new}}

Transactions are entered using the entry form (section \vref{transactions-form}), located below the list of transactions.
%La saisie d'une opération se fait dans le formulaire de saisie (section \vref{transactions-form}), situé en-dessous de la liste des opérations.


\subsection{Entering information\label{transactions-new-entry}}
%\subsection{Saisie des informations\label{transactions-new-entry}}

Enter the relevant details about your transaction in the various fields of the form. Some fields can be filled in either using the keyboard or via a drop-down list. If you use the keyboard, your entry will be automatically saved in the drop-down list for that field. If you wish, you can always remove this selection from the list later: refer to the various sections dealing with these lists (\menus{Payee}, \menus{Categories}, \menus{Budgetary lines} and the various tabs in the \menus{Resources} section of the \menus{Edit - Preferences} menu.
%Saisissez les différentes informations concernant votre opération dans les différents champs du formulaire. Certains champs peuvent être renseignés soit avec le clavier, soit avec une liste déroulante. Si vous utilisez le clavier, cette saisie sera automatiquement enregistrée dans la liste déroulante de ce champ. Si vous le voulez, vous pourrez toujours supprimer plus tard ce choix dans la liste: reportez-vous aux différents chapitres traitant de ces listes (\menus{Tiers}, \menus{Catégories}, \menus{Imputations budgétaires} et les différents onglets de la section \menus{Ressources} du menu \menus{Éditer - Préférences}.

Each entry field must be filled in with the appropriate data type (date, number, text); otherwise, the incorrect data will turn \textcolor{red}{red}.
%Chaque champ de saisie doit être renseigné par le type de donnée qui lui convient (date, nombre, texte), sinon la donnée erronée deviendra \textcolor{red}{rouge}.

Once the last field in the form has been filled in and submitted, the transaction appears in the list of transactions, and the active line, highlighted in \fboxsep=2pt{\colorbox[RGB]{245,155,155}{dark pink background}}, moves to the last empty line in the list.
%Une fois le dernier champ du formulaire saisi et validé, l'opération s'affiche dans la liste des opérations, et la ligne active, sur \fboxsep=3pt{\colorbox[RGB]{245,155,155}{fond rose foncé}}, passe à la dernière ligne vide de la liste.

% espace pour changement de thème
\vspacepdf{3mm}
The following subsections describe the various input methods available.
%Les sous-sections suivantes donnent les différentes méthodes de saisie disponibles.


\subsection{Main keyboard commands\label{transactions-new-keyboard}}
%\subsection{Commandes principales au clavier\label{transactions-new-keyboard}}

The \keys{Tab \tab} key allows you to navigate within the form.
%La touche \keys{Tab \tab} permet le déplacement dans le formulaire.

The \keys{Enter \return} key can be used either to move around the entry form or to submit the entry; this can be configured in the \menus{Edit - Preferences - Transaction form} menu (see the \vref{setup-form-behaviour} section, \menus{Form behavior}). This setting applies to both the \keys{Enter \return} key on the alphabetic keyboard and the one on the numeric keypad.
%La touche \keys{Entrée \return} permet soit de se déplacer dans le formulaire de saisie, soit de valider la saisie; cela est configurable dans le menu \menus{Éditer - Préférences - Formulaire des opérations} (voir la section \vref{setup-form-behaviour}, \menus{Comportement du formulaire}). Cette configuration affecte à la fois la touche \keys{Entrée \return} du clavier alphabétique et celle du pavé numérique.

The \keys{Esc} key cancels the current entry.
%La touche \keys{Échap} permet d'annuler la saisie en cours.


\subsection{Entering the date using the keyboard\label{transactions-new-dates}}
%\subsection{Saisie de la date au clavier\label{transactions-new-dates}}

The current date or the last date entered is displayed automatically when the form is opened.
%La date du jour ou la dernière date entrée s'affiche automatiquement à l'ouverture du formulaire.

%espace pour changement de thème
\vspacepdf{3mm}
You can enter the date in one of the following formats:
%Vous pouvez saisir la date avec l'un de ces formats:

\begin{itemize}
	\item mm/dd/yy (with a separator between the month, day and year);
	%\item jj/mm/aa (avec délimiteur entre le jour, mois et année);
	\item mmddyy or mmdd (without delimiters); in this case, you must use two digits for the month, day and year;
	%\item jjmmaa ou jjmm (sans délimiteur); dans ce cas, vous devrez bien utiliser deux chiffres pour les jour, mois et année;
	\item m/d/y if the month, day or year consists of a single digit (e.g. 1/1/1 for January 1st, 2001).
	%\item j/m/a si le jour, le mois ou l'année ne comportent qu'un seul chiffre (par ex.: 1/1/1 pour le premier janvier 2001).
\end{itemize}

Furthermore, there is no need to enter the year if it is the same as the current year, or the month if it is the current month, as Grisbi will fill these in automatically.
%De plus, il n'est pas nécessaire de saisir l'année si celle-ci est identique à l'année en cours, ni le mois si celui-ci est le mois en cours, car Grisbi complétera automatiquement.

%espace pour changement de thème
\vspacepdf{3mm}
Finally, you can increase or decrease the displayed date:
%Enfin, vous pouvez incrémenter ou décrémenter la date affichée:

\begin{itemize}
	\item per day using the \keys{{+}} or \keys{{-}} keys;
	%\item par jour avec les touches \keys{{+}} ou \keys{{-}};
	\item per week using the \keys{Shift \shift+{+}} or \keys{Shift \shift+{-}} keys;
	%\item par semaine avec les touches \keys{Maj \shift+{+}} ou \keys{Maj \shift+{-}};
	\item per month using the \keys{PageUp} or \keys{PageDown} keys;
	%\item par mois avec les touches\keys{Pg.Préc} ou \keys{Pg.Suiv};
	\item per year using the key combinations \keys{Shift \shift+PageUp} or \keys{Shift \shift+PageDown}.
	%\item par année avec les combinaisons de touches \keys{Maj \shift+Pg.Préc} ou \keys{Maj \shift+Pg.Suiv}.
\end{itemize}


\subsection{Entering a date in the calendar\label{transactions-new-calendar}}
%\subsection{Saisie de la date au calendrier\label{transactions-new-calendar}}

Double-clicking on the date, or pressing the \keys{Ctrl+Enter \return} key combination, brings up a calendar from which you can select the date.
%Un double-clic sur la date, ou la combinaison de touches \keys{Ctrl+Entrée \return}, affiche un calendrier permettant de sélectionner la date.

%espace pour changement de thème
\vspacepdf{3mm}
You can then select a new date using the mouse or keyboard:
%Vous pouvez alors sélectionner une nouvelle date avec la souris ou le clavier:
% espace avant image 5mm
\vspacepdf{3mm}

\noindent
\begin{minipage}{.65\linewidth}
	\begin{itemize}[rightmargin=.6cm]
		\item the \keys{{+}} or \keys{{-}} keys, as well as the \keys{\arrowkeyleft} or \keys{\arrowkeyright} keys, allow you to move to the previous or next day;
		%\item les touches \keys{{+}} ou \keys{{-}}, de même que \keys{\arrowkeyleft} ou \keys{\arrowkeyright}, permettent de passer au jour précédent ou suivant;
		\item the \keys{\arrowkeyup} or \keys{\arrowkeydown} keys allow you to move to the previous or next week;
		%\item les touches \keys{\arrowkeyup} ou \keys{\arrowkeydown} permettent de passer à la semaine précédente ou suivante;
		\item the \keys{PageUp} or \keys{PageDown} keys allow you to move to the previous or next month;
		%\item les touches \keys{Pg.préc} ou \keys{Pg.Suiv} permettent de passer au mois précédent ou suivant;
		\item the \keys{Home} or \keys{End} keys allow you to move to the first or last day of the month;
		%\item les touches \keys{Début} ou \keys{Fin} permettent de passer au premier jour ou au dernier jour du mois;
		\item clicking the \menus{Cancel} button or pressing the \keys{Esc} key closes the calendar without changing the date;
		%\item un clic sur le bouton \menus{Annuler} ou l'appui sur la touche \keys{Échap} ferme le calendrier sans modifier la date;
		\item double-clicking or pressing the \keys{Enter \return} key confirms the selected date.
		%\item un double-clic ou l'appui sur la touche \keys{Entrée \return} valide la date sélectionnée.
	\end{itemize}
\end{minipage}
\hspace{10pt}
\begin{minipage}{.35\linewidth}
	%\begin{wrapfigure}{R}{0.4\textwidth}%50mm
	%\vspace{-\intextsep}				% space above the floating (minus intextsep=separation between float and text in text)
	\vspace{-5pt}
	\centering							% centering the floating figure in the "wrap"
	%\raggedleft
	\includegraphics[width=1\textwidth]{image/screenshot/transactions_calendar}
	\vspace{-15pt}						% space between the floating and the caption below the floating
	\captionsetup{%						% for a change of options of the caption package
		format=plain,					% to avoid error message with labelsep option below, default=hang
		type=figure,%
		name=Fig.,%						% rename label caption, default="Figure" (or "Table")
		justification=centering,		% centring of text and caption label 
		labelsep=newline}				% define the separator between the label and the caption text
	\caption{Calendar}
	%\caption{Calendrier}		% \caption is mandatory to reference a figure in lof
	\vspace{-10pt}						% space below the caption
	\label{transactions-calendar-img}
	%\end{wrapfigure}
\end{minipage}


\subsection{Drop-down lists\label{transactions-new-lists}}
%\subsection{Listes déroulantes\label{transactions-new-lists}}

Some fields have a drop-down list of predefined options. You can select an option using the mouse, or you can type the data into the entry field using the keyboard.
%Certains champs disposent d'une liste déroulante de choix déjà définis. Vous pouvez y sélectionner un libellé avec la souris, ou préférer saisir les données au clavier dans le champ de saisie.

You can open a drop-down list by clicking the small triangle to the right of the field with the mouse, or by pressing the \keys{PagegDown} or \keys{\arrowkeydown} keys on the keyboard; to navigate through the list, you can use:
%L'ouverture d'une liste déroulante peut se faire avec la souris en cliquant sur le petit triangle à droite du champ, ou avec les touches du clavier \keys{Pg.Suiv} ou \keys{\arrowkeydown}, le déplacement dans cette liste peut se faire avec:
\begin{itemize}
	\item the mouse and the scrollbar to the right of the list,
	%\item la souris et l'ascenseur à droite de la liste,
	\item the mouse wheel,
	%\item la molette de la souris,
	\item the \keys{PageUp}, \keys{PageDown}, \keys{\arrowkeyup} or \keys{\arrowkeydown} keys,
	%\item les touches \keys{Pg.Préc}, \keys{Pg.Suiv}, \keys{\arrowkeyup} ou \keys{\arrowkeydown},
\end{itemize}
	
%et la validation d'un choix à l'intérieur d'une liste déroulante se fait avec la touche \keys{Entrée \return}.

\vspacepdf{3mm}
\Note{}: depending on the theme of your desktop environment or the window manager you are using, these triangles can be replaced with other characters such as +, -, >, <, etc.
%\Note{}: ces triangles peuvent être remplacés, en fonction du thème de l'environnement de bureau ou du gestionnaire de fenêtres que vous utilisez, par d'autres caractères tels que +, -, >, <, etc.

%espace pour changement de thème
\vspacepdf{3mm}
Grisbi will suggest automatically completing the word you have entered as soon as it has recognised enough characters. The selected \indexword{auto-complete}\index{entry field!auto-complete} appears on \fboxsep=3pt{\colorbox[RGB]{61,174,233}{blue background}}. To confirm your choice, press the \keys{Tab \tab} key or the \keys{Enter \return} key, depending on your configuration settings; otherwise, continue typing.
%Grisbi vous proposera de compléter automatiquement le mot saisi dès qu'il en aura reconnu suffisamment de caractères. L'\indexword{auto-complétion}\index{champ de saisie!auto-complétion} sélectionnée apparaît sur \fboxsep=3pt{\colorbox[RGB]{61,174,233}{fond bleu}}. Pour valider votre choix, appuyez sur la touche \keys{Tab \tab} ou la touche \keys{Entrée \return}, selon votre choix de configuration, sinon continuez la saisie.

The drop-down lists for currencies, financial years and methods of payment can be opened using the \keys{Space} key or a \texttt{left-click} of the mouse; you can navigate through them using the \keys{\arrowkeyup} and \keys{\arrowkeydown} keys, and confirm your selection using the \keys{Space} key, the \keys{Enter \return} key or a \texttt{left-click} of the mouse.
%Les listes déroulantes des devises, exercices et modes de règlement s'ouvrent avec la touche \keys{Espace} ou un clic gauche de la souris; on s'y déplace avec les touches \keys{\arrowkeyup} et \keys{\arrowkeydown} et l'on valide avec la touche \keys{Espace}, la touche \keys{Entrée \return} ou un clic gauche de la souris.


\subsection{Entering formulas\label{transactions-new-numbers}}
%\subsection{Saisie de formules\label{transactions-new-numbers}}

In the entry form, you can enter numbers, with or without decimals, into fields such as \menus{Debit} and \menus{Credit} using the keyboard.
%Dans le formulaire de saisie, vous pouvez saisir les nombres, avec ou sans décimales, dans les champs tels que \menus{Débit} et \menus{Crédit}, à l'aide du clavier.

You can also enter simple \indexword{arithmetic expressions}\index{entry field!arithmetic expression}\index{arithmetic expression!entry field} using the four basic arithmetic operators: addition, subtraction, multiplication and division. This replaces the \indexword{calculator}\index{calculator} on your desk.
%Vous pouvez aussi y saisir des \indexword{formules arithmétiques}\index{champ de saisie!formule arithmétique}\index{formule arithmétique!champ de saisie} simples, à l'aide des quatre opérateurs d'addition, soustraction, multiplication et division. Cela remplace la \indexword{calculatrice}\index{calculatrice} sur votre bureau.

To enter an arithmetic formula, type it as if it were text, for example \enquote{3.40-2.10}, then continue entering the other fields of your transaction as normal; if it is entered incorrectly, it will be displayed in \textcolor{red}{red}. When you click the \menus{Validate} button and for a new transaction, the result \enquote{1.30} appears in the list of transactions; otherwise, the result also appears in the entry field. And if the formula is still incorrect, the message \enquote{\textcolor{red}{\#\#\#\textsf{ERR}\#\#\#}}  will appear instead and you will not be able to validate the financial transaction.
%Pour saisir une formule arithmétique, saisissez-la comme si c'était un texte, par exemple \frquote{3,40-2,10}, puis continuez normalement la saisie des autres champs de votre opération; si elle est mal écrite, elle s'affiche en \textcolor{red}{rouge}. Lorsque vous cliquez sur le bouton \menus{Valider} et pour une nouvelle opération, le résultat \frquote{1,30} s'affiche dans la liste des opérations, sinon le résultat s'affiche, en plus, dans le champ de saisie. Et si la formule est toujours mal écrite, le message \frquote{\textcolor{red}{\#\#\#\textsf{ERR}\#\#\#}}  s'affiche à la place et vous ne pourrez pas valider l'opération financière.

\vspacepdf{2mm}
\Attention{}: please use only the period as a decimal point and the comma as a thousands separator; otherwise, this will cause calculation errors (e.g., "1,2+2.4=14.40"). For your information, with Grisbi in other languages (French, German, etc.), the period and comma are interchangeable as decimal separators.
% FAUX pour la version anglaise : la virgule sert de séparateur de milliers et fait un "mauvais" mélange avec le point comme séparateur décimal.
%Cette calculatrice accepte le mélange des points et des virgules comme séparateur décimal, par exemple \frquote{1.2+2,3}, et Grisbi effectuera le calcul.
\vspacepdf{2mm}

This calculator accepts the multiplication or division, but in a manner exclusive of addition or subtraction; that is to say that we can mix additions and subtractions in an entry but not a multiplication (or division) and any other operator: we cannot put parentheses to indicate an order of priority of operation.
%Cette calculatrice accepte la multiplication ou la division, mais de manière exclusive de l'addition ou de la soustraction; c'est-à-dire qu'on peut mélanger additions et soustractions dans une saisie mais pas une multiplication (ou une division) et un quelconque autre opérateur: on ne peut pas mettre de parenthèses pour indiquer un ordre de priorité d'opération.

This feature can be very useful, particularly for adjusting the amounts of a split transaction in the event of a discrepancy.
%Cette fonctionnalité peut être très utile, notamment pour réajuster les montants d'une opération ventilée en cas d'écart.


\subsection{Depositing cheques or cash\label{transactions-new-cheque}}
%\subsection{Remise de chèques ou d'espèces\label{transactions-new-cheque}}

To record a \indexword{cheque deposit}\index{deposit!cheque}, first deposit the cheques into your bank account at the bank, which will give you a deposit slip, then enter the deposit transaction as a credit to your bank account.
%Pour enregistrer une \indexword{remise de chèques}\index{remise!chèques}, vous les déposez d'abord sur votre compte bancaire à la banque qui vous remet un bordereau de remise, puis vous saisissez l'opération de remise, en crédit sur votre compte bancaire.

To record a \indexword{cash deposit}\index{deposit!cash}, you must first deposit the cash either into your bank account at the bank, which will issue you with a deposit slip, or into your cash box, and then enter the deposit transaction as a credit to your cash account.
%Pour enregistrer une \indexword{remise d'espèces}\index{remise!espèces}, vous les déposez d'abord soit sur votre compte bancaire à la banque qui vous remet un bordereau de remise, soit dans votre porte-monnaie, puis vous saisissez l'opération de remise, en crédit sur votre compte de caisse.

% espace pour changement de thème
\vspacepdf{3mm}
If you receive a lot of cheques (for example, because you run a charity and your members pay their membership fees by cheque), you have two options:
%Si vous recevez beaucoup de chèques (par exemple parce que vous gérez une association et que vos adhérents payent leur cotisation par chèque), vous avez deux méthodes de gestion:

\begin{itemize}
	\item either you record each cheque as a separate transaction in your bank account, and then bundle these cheques together to deposit them at the bank; your deposit slip will then show only the total amount; when reconciling your bank account, your bank statement will show only this total amount, but not the amount of each cheque entered into your bank account; as a result, reconciliation becomes rather difficult (you should at least have made a list of these cheques beforehand);
	%\item soit vous enregistrez individuellement une opération pour chaque chèque dans votre compte bancaire, puis vous groupez ces chèques pour les déposer à la banque; votre bordereau de remise ne contient alors que leur montant global; lors du rapprochement bancaire, votre relevé bancaire ne mentionne que ce montant global, mais non le montant de chaque chèque saisi dans votre compte bancaire; de ce fait, le pointage est rendu peu facile (vous devriez au moins avoir fait auparavant une liste de ces chèques);
	\item either you record a single transaction for all your cheques, in which case you will have lost the third-party details for each cheque, and this information will not appear in the \menus{Payee} tab, even if you try to split this transaction.
	%\item soit vous enregistrez une seule opération pour tous vos chèques, vous aurez alors perdu l'information du tiers de chaque chèque, et cette information n'apparaîtra pas dans l'onglet \menus{Tiers}, même si vous essayez de ventiler cette opération.
\end{itemize} 

% espace pour changement de thème
\vspacepdf{3mm}
If none of these methods suit you, a better solution is to use a \emph{suspense account}. The procedure for doing so is explained in detail in the section \vref{accounts-type-bank-waiting-remittance}, \menus{Delivery of cheques or cash}.
%Si aucune de ces méthodes ne vous convient, une meilleure solution consiste à utiliser un \emph{compte d'attente}. Pour cela, la méthode est explicitée en détails au paragraphe \vref{accounts-type-bank-waiting-remittance}, \menus{Remise de chèques ou d'espèces}.


\subsection{Transfers\label{transactions-new-transfer}}
%\subsection{Virements\label{transactions-new-transfer}}

The term \indexword{transfer}\index{transaction!transfer} covers two different concepts: external transfers and internal transfers. Here are the recommended entry methods for each case.
%Le terme \indexword{virement}\index{opération!virement} recouvre deux notions différentes: les virements externes et les virements internes. Voici dans chaque cas les méthodes de saisie conseillées.


\subsubsection{External transfers\label{transactions-new-transfer-external}}
%\subsubsection{Virements externes\label{transactions-new-transfer-external}}

\indexword{External transfers}\index{external transfer} are transfers made by your bank, or by you if you place an order online, to or from your account. They can be incoming (e.g. if your employer instructs your bank to transfer your salary to you) or outgoing (e.g. if you instruct your bank to make a SWIFT transfer to pay a bill abroad).
%Les \indexword{virements externes}\index{virement!externe} sont des virements que votre banque, ou vous-même si vous passez un ordre par Internet, effectue vers ou à partir de votre compte. Ils peuvent être entrants (par ex. si votre employeur donne l'ordre à votre banque de vous virer le montant de votre salaire) ou sortants (par ex. si vous donnez l'ordre à votre banque de faire un virement Swift pour payer une facture à l'étranger).

This is a type of transaction similar to a cheque, and you can enter it by selecting \menus{Payment of method}: \menus{Transfer} in the entry form.
%Il s'agit d'un type d'opération assimilable à un chèque, et vous pourrez la saisir en sélectionnant le \menus{Moyen de paiement}: \menus{Virement} dans le formulaire de saisie.


\subsubsection{Internal transfers\label{transactions-new-transfer-internal}}
%\subsubsection{Virements internes\label{transactions-new-transfer-internal}}

The \indexword{internal transfers}\index{transfer!internal} are transfers between your various Grisbi accounts, which do not necessarily have an equivalent in the bank; here are some examples:
%Les \indexword{virements internes}\index{virement!interne} sont des transferts entre vos différents comptes de Grisbi, et qui n'ont pas forcément d'équivalent en banque; voici quelques exemples:

\begin{itemize}
	\item if you use a cash account, when you make a cash withdrawal at the bank, you enter a transfer to that cash account in Grisbi, but the transaction type is a \menus{Withdrawal};
	%\item si vous utilisez un compte de caisse, lorsque vous faites un retrait d'espèces à la banque, vous saisissez dans Grisbi un virement vers ce compte de caisse, mais le type d'opération est \menus{Retrait};
	\item if you use an advance account, when you issue or receive a refund check, you enter an internal transfer, even though the type of transaction is \menus{Check};
	%\item si vous utilisez un compte d'avances, lorsque vous établissez ou recevez un chèque de remboursement, vous saisissez un virement interne, bien que le type d'opération soit \menus{Chèque};
	\item if you use a liability account to manage a loan, you enter the monthly direct debit made by your bank as an internal transfer from your bank account to your liability account, as a \menus{Direct debit};
	%\item si vous utilisez un compte de passif pour gérer un emprunt, vous saisissez le prélèvement mensuel effectué par votre banque comme un virement interne de votre compte bancaire vers votre compte de passif, en tant que type \menus{Prélèvement};
	\item if you use a bank account for a deferred debit card, you enter the monthly withdrawal by your bank as a transfer from your current account to this card account.
	%\item si vous utilisez un compte bancaire pour une carte bancaire à débit différé, vous saisissez le prélèvement mensuel par votre banque comme un virement de votre compte courant vers ce compte de carte.
\end{itemize}

However, if you do not use any of these accounts and your account file contains only a current account and a savings account, then internal transfers will be transfers from one account to the other. These will indeed be transactions of the \menus{Transfer} type, and you should record them in the \menus{Transfer: Savings account} or \menus{Transfer: Current account} category, as appropriate.
%Par contre, si vous n'utilisez aucun de ces comptes et que votre fichier de comptes contient uniquement un compte courant et un compte d'épargne, alors les virements internes correspondent à des virements d'un compte vers l'autre. Ce seront bien des opérations de type \menus{Virement}, et vous les enregistrerez dans la catégorie \menus{Virement: Compte d'épargne} ou \menus{Virement: Compte courant}, selon le cas.

You can make transfers between accounts by selecting the category \emph{Transfer: the\_account\_to\_be\_charged}. Depending on whether you enter the amount under \menus{Debit} or \menus{Credit}, this will be a transfer \emph{to the account} or a transfer \emph{from the account}. Of course, both accounts must be in the same file created by Grisbi.
%Vous pouvez réaliser des virements entre comptes en sélectionnant la catégorie \emph{Virement: le\_compte\_à\_imputer}. Selon que vous saisissez son montant en \menus{Débit} ou en \menus{Crédit}, ce sera un virement \emph{vers le compte} ou un virement \emph{à partir du compte}. Il est bien sûr nécessaire que les deux comptes se trouvent dans le même fichier créé par Grisbi.

% espace pour changement de thème
\vspacepdf{2mm}
If you make a \indexword{transfer between two accounts in different currencies}\index{currency!different}\index{currency!transfer}\index{transfer!currency}, Grisbi will automatically carry out the conversion, and the debit and credit will clearly show the currencies; for example, if, in your yen account, you make a transfer transaction in credit from an pound sterling account, the transfer will appear in yen in the \menus{Credit} column, and its yen balance will be correct. In the pound sterling account, the debit amount will appear in yen in the \menus{Debit} column, and in pound sterling in parentheses, in the \menus{Amount} (\vref{transactions-list-fields}) field in the \menus{Credit} column (by default), if the display is at least in \menus{Two lines view} (see notes below).
%Si vous faites un \indexword{virement entre deux comptes de devises différentes}\index{devise!différente}\index{devise!virement}\index{virement!devise}, Grisbi effectuera automatiquement la conversion, et le débit et le crédit feront apparaître clairement les devises; par exemple, si, dans votre compte en yens, vous faites une opération de virement en crédit provenant d'un compte en euros, le virement apparaîtra en yens dans la colonne \menus{Crédit}, et son solde en yens sera correct. Dans le compte en euros, le montant du débit apparaîtra en yens dans la colonne \menus{Débit}, et en euro entre parenthèses, dans le champ \menus{Montant} (\vref{transactions-list-fields}) dans la colonne \menus{Crédit} (par defaut), si l'affichage est au minimum en \menus{Mode \frquote{deux lignes}} (voir les notes ci-dessous).

\vspacepdf{2mm}
\noindent
\Note{}: to display the \menus{Amount} field, you will need to configure the transaction list by:
%\Note{}: pour que le champ \menus{Montant} s'affiche, vous devrez configurer la liste des opérations en: 
\begin{itemize}
	\item checking/adding the \menus{Amount} field (see the section \vref{setup-operations-cells}, \menus{Transaction list cells} or the section \vref{transactions-list-fields-manage}, \menus{Management of information fields}),
	%\item vérifiant/ajoutant le champ \menus{Montant} (voir la section \vref{setup-operations-cells}, \menus{Cellules de la liste des opérations} ou la section \vref{transactions-list-fields-manage}, \menus{Gestion des champs d'information}), 
	\item adjusting the number of lines displayed in the transaction list (see the section \vref{transactions-functions} \menus{Toolbar} or the section \vref{home-menus-display} \menus{View Menu}).
	%\item ajustant le nombre de lignes affichées dans la liste des opérations (voir la section \vref{transactions-functions} \menus{Barre d'outils} ou la section \vref{home-menus-display} \menus{Menu Affichage}).
\end{itemize}

\vspacepdf{2mm}
\Note{}: if you display an existing transaction in the entry form, you should configure the form’s content by adding the \menus{Currency} field to view the currency associated with the debit or credit transaction (see the \vref{setup-form-content} section, \menus{Content}).
%\Note{}: si vous affichez une opération existante dans le formulaire de saisie, vous devrait configurer le contenu du formulaire de saisie en y ajoutant le champ \menus{Devise} pour visualiser la devise associée à l'opération en débit ou en crédit (voir la section \vref{setup-form-content}, \menus{Contenu}).


\subsubsection{Contra-transaction of a transfer\label{transactions-new-transfer-associatedOperation}}
%\subsubsection{Contre-opération d'un virement\label{transactions-new-transfer-associatedOperation}}

When you enter a transaction of the type \enquote{internal transfer} - from one account to another, or from one account to another - Grisbi automatically creates the \indexword{contra-transaction}\index{contra-transaction}\index{transaction!contra-transaction}\index{transfer!contra-transaction} in the other account, i.e. the corresponding transfer in the recipient account or the source account.
%Lorsque vous saisissez une opération de type \frquote{virement interne}, d'un compte vers un autre, ou sur un compte à partir d'un autre, Grisbi crée automatiquement la \indexword{contre-opération}\index{contre-opération}\index{opération!contre-opération}\index{virement!contre-opération} dans l'autre compte, c'est-à-dire le virement correspondant dans le compte destinataire ou dans le compte origine.

If the transfer is in a foreign currency, it will automatically create the contra-transaction entry in the currency of the destination account if the exchange rate is fixed; otherwise, it will ask you to confirm this before posting the entry.
%S'il s'agit d'un virement en devises, il créera aussi la contre-opération dans la devise du compte de destination de façon automatique si le taux de change de la devise est fixe, sinon il vous le demandera avant de passer l'écriture.

% espace pour changement de thème
\vspacepdf{3mm}
To \indexword{view the contra-transaction}\index{contra-transaction!view} for a transfer, \texttt{right-click} on its line in one of the two accounts, and in the context menu, select \menus{Displays the contra-transaction}: the other account is displayed and the transaction is selected (on \fboxsep=3pt{\colorbox[RGB]{245,155,155}{dark pink background}}). If you repeat exactly the same action in that same account, the original account and its transfer are displayed again.
%Pour \indexword{afficher la contre-opération}\index{contre-opération!afficher} d'un virement, cliquez-droit sur sa ligne dans un des deux comptes, et dans le menu contextuel, sélectionnez \menus{Afficher la contre-opération}: l'autre compte s'affiche et l'opération est sélectionnée (sur \fboxsep=3pt{\colorbox[RGB]{245,155,155}{fond rose foncé}}). Si vous refaites exactement la même manipulation dans ce même compte, le compte d'origine et son virement se ré-affichent.


\subsection{Advances: consent and receipt\label{transactions-new-advance}}
%\subsection{Avances: consentement et réception\label{transactions-new-advance}}

You may need to \indexword{consent to or receipt an advance}\index{advance!consent}\index{advance!receipt} of funds: this could be, for example, an advance on your salary or a purchase you make on behalf of a friend.
%Vous pouvez être amené à \indexword{consentir ou recevoir une avance}\index{avance!consentement}\index{avance!réception} de fonds: il peut s'agir, par exemple, d'une avance sur votre salaire ou d'un achat que vous faites pour le compte d'un de vos amis.

%espace pour changement de thème
\vspacepdf{2mm}
There are two ways to handle this:
%Vous pouvez gérer cela de deux manières:
\begin{itemize}
	\item the simplest approach is to create a dedicated category and assign debit and credit transactions to it; once all advances have been repaid, the balance of this category, as shown in the list of categories, should be zero (see the section \vref{categories-list}, \menus{List of categories and sub-categories});
	%\item la plus simple est de créer une catégorie dédiée, et d'y affecter les opérations de débit et de crédit; quand toutes les avances sont remboursées, le solde de cette catégorie, affiché dans la liste des catégories, doit être égal à zéro (voir la section \vref{categories-list}, \menus{Liste des catégories et des sous-catégories}); 
	\item a more useful approach is to create an advance account, which provides a more immediate overview of the status of outstanding advances: this is described in detail in the section \vref{accounts-type-bank-advance}, \menus{Advance Account}.
	%\item une manière plus intéressante est de créer un compte d'avances, qui permet une vision plus immédiate de l'état des avances en cours: elle est décrite en détails dans le paragraphe \vref{accounts-type-bank-advance}, \menus{Compte d'avances}.
\end{itemize}


\subsection{Automatic form completion\label{transactions-new-autofill}}%\label{transactions-new-recall}
%\subsection{Remplissage automatique du formulaire\label{transactions-new-autofill}}%\label{transactions-new-recall}

The \indexword{automatic completion}\index{form!automatic completion}\index{transaction!automatic completion} of a transaction, if you have enabled the \menus{Automatic filling transactions from payee} option (under the \menus{Edit - Preferences - Transaction form - Completion}), allows you to enter a transaction assigned to a payee already known to your Grisbi file much more quickly: when entering the payee, as soon as you leave the input field, the other fields in the form are automatically filled in to match the last transaction entered for that payee. It is up to you to modify them if necessary before confirming the transaction.
%Le \indexword{remplissage automatique}\index{formulaire!remplissage automatique}\index{opération!remplissage automatique} d'une opération, si vous avez activé l'option \menus{Remplissage automatique des opérations à partir du tiers} (du menu \menus{Éditer - Préférences - Formulaire des opérations - Aide à la saisie}), vous permet de saisir beaucoup plus rapidement une opération affectée à un tiers déjà connu par votre fichier Grisbi: lors de la saisie du tiers, dès que vous sortez de la zone de saisie, les autres champs du formulaire sont automatiquement complétés à l'identique de la dernière opération saisie pour ce tiers. À vous de les modifier si nécessaire avant de valider l'opération.

You can fine-tune this behaviour in the \menus{Preferences} menu (see the section \vref{setup-form-input}, \menus{Form completion}).
%Vous pouvez ajuster finement ce comportement dans les \menus{Préférences} (voir la section \vref{setup-form-input}, \menus{Aide à la saisie du formulaire}).

% espace avant Attention ou Note: 5 mm
\vspacepdf{3mm}
\Note{}: if you do not change any fields and confirm the transaction, you will simply end up with a duplicate transaction, but one dated today.
%\Note{}: si vous ne modifiez aucun champ et que vous validez l'opération, vous obtiendrez une opération purement et simplement dupliquée, mais qui portera la date du jour courant.


\section{Entering a currency transaction\label{transactions-currencies}}
%\section{Saisie d'une opération en devises\label{transactions-currencies}}

Grisbi can handle transactions entered in a \indexword{different currency}\index{currency!different} from that of the account, with or without exchange fees.
%Grisbi peut gérer des opérations saisies avec une \indexword{devise différente}\index{devise!différente} de celle du compte, avec ou sans frais de change.

% espace pour changement de thème
\vspacepdf{3mm}
To use a different currency, you must first carry out the following settings:
%Pour utiliser une devise différente, vous devez d'abord procéder aux configurations suivantes:
\begin{enumerate}
	\item create the currency, if it has not already been created, via the \menus{Edit - Preferences - Resources} menu (see the \vref{setup-resources-currencies} section, \menus{Currencies});
	%\item créez la devise, si elle ne l'est pas déjà, dans le menu \menus{Éditer - Préférences- Ressources} (voir la section \vref{setup-resources-currencies}, \menus{Devises});
	\item in the following menu, create the link(s) you need between the new currency and the others already saved (see the section \vref{setup-resources-rate}, \menus{Currencies links});
	%\item dans le menu suivant, créez le ou les liens dont vous avez besoin entre la nouvelle devise et les autres déjà enregistrées (voir la section \vref{setup-resources-rate}, \menus{Liens entre devises});
	\item configure the display of the \menus{Currency} field in the transaction form (see the \vref{setup-form-content} section, \menus{Content});
	%\item configurez l'affichage du champ \menus{Devise} dans le formulaire de saisie des opérations (voir la section \vref{setup-form-content}, \menus{Contenu});
	% saut de ligne pour alignement de la note sur l'item
	 
	\Note{}: the \menus{Currency} and \menus{Change} fields are linked in the data entry form: configuring the display of the \menus{Currency} field in the form automatically positions the \menus{Change} field in the same place, and vice versa; you can then move them to the desired cells.
	%\Note{}: les champs \menus{Devise} et \menus{Change} sont liés dans le formulaire de saisie: la configuration de l'affichage du champ \menus{Devise} dans le formulaire y positionne automatiquement le champ \menus{Change}, et inversement; vous pouvez ensuite les déplacer dans les cellules voulues.
	\item configure the display of the \menus{Amount} field in the transaction list (see the section \vref{setup-operations-cells}, \menus{Transactions list cells} or the section \vref{transactions-list-fields-manage}, \menus{Management of information fields}).
	%\item configurez l'affichage du champ \menus{Montant} dans la liste des opérations (voir la section \vref{setup-operations-cells}, \menus{Cellules de la liste des opérations} ou la section \vref{transactions-list-fields-manage}, \menus{Gestion des champs d'information}).
\end{enumerate}

% espace pour changement de thème
\vspacepdf{3mm}
When entering the \indexword{currency transaction}\index{transaction!currency}\index{currency!transaction} in the entry form, select the currency in which the transaction is to be recorded from the drop-down list. Several questions may arise:
%Lors de la saisie de l'\indexword{opération en devise}\index{opération!devise}\index{devise!opération} dans le formulaire de saisie, choisissez dans la liste déroulante la devise dans laquelle l'opération doit être enregistrée. Plusieurs questions peuvent se poser:

% image centrée 
\begin{figure}[htbp]
\begin{center}
\includegraphics[width=0.95\textwidth]{image/screenshot/transactions_changeRate}
\end{center}
\caption{Fees and exchange rates}
%\caption{Frais et taux de change}
\label{transactions-changeRate-img}
\end{figure}
% image centrée 

\begin{itemize}
	\item if you have not created a link (exchange rate) for this currency, Grisbi displays the fees and exchange rates window \refimage{transactions-changeRate-img}: enter the exchange rate there, confirm the window, then confirm the transaction: the amount in the currency is displayed in the \menus{Debit} field and the amount in the account’s currency in the \menus{Amount} field;
	%\item si vous n'avez pas créé de lien (taux de change) pour cette devise, Grisbi affiche la fenêtre de frais et taux de change \refimage{transactions-changeRate-img}: saisissez-y le taux de change, validez la fenêtre puis validez l'opération: le montant en devise s'affiche dans le champ \menus{Débit} et celui dans la devise du compte dans le champ \menus{Montant};
	% saut de ligne pour alignement de la note sur l'item
	 
	\Note{}: you can reverse the direction of the conversion (one pound sterling is worth X Currency, or one Currency is worth Y pounds sterling) by swapping the names of the currencies in the two drop-down lists; this allows you to enter the exchange rate exactly as it has been provided to you; this rate can be adjusted at any time after the transaction has been created.
	%\Note{}: vous pouvez inverser le sens de la conversion (un euro vaut X Devise, ou bien une Devise vaut Y euros) en inversant le nom des devises dans les deux listes déroulantes; vous pouvez ainsi entrer exactement le taux de change tel qu'il vous est communiqué; ce taux peut être ajusté à tout moment après la création de l'opération.
	\item if you wish to change the exchange rate, open the transaction for editing, click the \menus{Change} button in the input form, tick the \menus{Change the link} checkbox if it is not already ticked, then change the rate, and save the window and the transaction;
	%\item si vous voulez modifier le taux de change, ouvrez l'opération en modification, cliquez sur le bouton \menus{Change} dans le formulaire de saisie, cochez la case de libellé \menus{Modifier le taux de change} si elle ne l'est pas, puis modifiez-le, et validez la fenêtre puis l'opération;
	\item if you have any \indexword{exchange fees}\index{exchange fees}, open the transaction in edit mode, click the \menus{Change} button in the entry form, enter the fees in the account’s currency in the \menus{Exchange fees} field, and confirm the window and then the transaction; These fees are then added to the contents of the \menus{Amount} field; however, you can instead enter them in a separate transaction, using a dedicated category, for example \emph{Financial expenses: Bank charges}: this allows you to record them separately;
	%\item si vous avez des \indexword{frais de change}\index{frais de change}, ouvrez l'opération en modification, cliquez sur le bouton \menus{Change} dans le formulaire de saisie, saisissez les frais en devise du compte dans le champ \menus{Frais de change}, et validez la fenêtre puis l'opération; ces frais sont alors ajoutés au contenu du champ \menus{Montant}; mais vous pouvez à la place les saisir dans une opération indépendante, en utilisant une catégorie dédiée, par exemple \emph{Frais financiers: Frais bancaires}: cela peut vous permettre de les comptabiliser à part;
	\item if you want the exchange rate to be fixed: see the section \vref{setup-resources-rate}, \menus{Currencies links}.
	%\item si vous voulez que le taux de change soit fixe: voir la section \vref{setup-resources-rate}, \menus{Liens entre devises}.
\end{itemize}


\section{Entering a transaction with a virtual third party\label{transactions-virtualThird}}
%\section{Saisie d'une opération avec tiers virtuel\label{transactions-virtualThird}}

Grisbi allows you to use virtual third parties: a virtual third party is a \emph{report} that represents a list of several third parties.
%Grisbi vous permet d'utiliser des tiers virtuels: un tiers virtuel est un \emph{état} qui représente une liste de plusieurs tiers.

When you enter a transaction with a virtual third party as the payee, Grisbi records, upon validation, an identical transaction (amount, category, budgetary line, payment method, etc.) for \emph{each} of the third parties represented by that virtual third party. For example, you can enter a membership fee request for 200 members of an association in a single go, which saves a considerable amount of time\ldots
%Lorsque vous saisissez une opération avec pour tiers un tiers virtuel, Grisbi enregistre, au moment de sa validation, une opération identique (montant, catégorie, imputation budgétaire, moyen de paiement etc.) pour \emph{chacun} des tiers représentés par ce tiers virtuel. Par exemple, vous pouvez saisir en une seule fois un appel à cotisation pour 200 adhérents d'une association, ce qui représente un gain de temps très appréciable\ldots

\vspacepdf{3mm}
To \indexword{enter a transaction with a virtual third party}\index{virtual third party!transaction entry}, proceed as follows:
%Pour \indexword{saisir une opération avec un tiers virtuel}\index{tiers virtuel!saisie d'opération}, procédez de la sorte:
\begin{itemize}
	\item enter this third party in the \menus{Payee} field of the entry form, in the format \enquote{Report : name\_of\_virtual\_third\_party};
	%\item saisissez ce tiers dans le champ \menus{Tiers} du formulaire de saisie, sous la forme \frquote{État : nom\_du\_tiers\_virtuel}.
	% saut de ligne pour alignement de la note sur l'item
	
	\Note{}: please ensure you follow the syntax correctly, with a space before AND after the \enquote{:}.
	%\Note{}: veuillez à bien respectez la syntaxe avec un espace avant ET après les \frquote{:}.
	\item create a virtual third party by creating a report of type \menus{Blank report} containing a list of payees: see the section \vref{reportscreation-display-general-virtualThird}, \menus{Virtual Third Party};
	%\item créez un tiers virtuel en créant un état de type \menus{État vierge} contenant une liste de tiers: voir le paragraphe \vref{reportscreation-display-general-virtualThird}, \menus{Tiers virtuel},
	\item modify the virtual third party as required by editing the report that created it: see the section \vref{reports-modify}, \menus{Modifying a report}.
	%\item modifiez le tiers virtuel à votre convenance en modifiant l'état qui l'a créé: voir la section \vref{reports-modify}, \menus{Modification d'un état}.
\end{itemize}

%espace pour changement de thème
\vspacepdf{2mm}
To delete a virtual third party, you have two options:
%Pour supprimer un tiers virtuel, vous avez deux possibilités:
\begin{itemize}
	\item uncheck the \menus{Consider the payees of this report as a multiple payee} box in the \menus{General} section of the report that defines it, then click on the \menus{Validate} button (see the section \vref{reports-modify},
	%\item décochez la case \menus{Considérer les tiers de cet état comme un tiers virtuel} dans les \menus{Généralités} de l'état qui le définit, puis validez (voir la section \vref{reports-modify}, \menus{Modification d'un état});
	\item delete the report that defines it (see the section \vref{reports-delete}, \menus{Deleting a report}).
	%\item supprimez l'état qui le définit (voir la section \vref{reports-delete}, \menus{Suppression d'un état}).
\end{itemize}


\section{Breakdown of a transaction\label{transactions-breakdown}}
%\section{Ventilation d'une opération\label{transactions-breakdown}}

\indexword{Break down a transaction}\index{break down!transaction} means to split the amount across several lines (for example, your supermarket purchases can be split between the categories \emph{Food : Groceries} and \emph{Personal care : Clothing}). You can also split a sum within the same category, but with different information in the \menus{Notes} field (for example, to record the details of a transaction).
%\indexword{Ventiler une opération}\index{ventiler!opération} signifie répartir son montant entre plusieurs lignes (par exemple, vos achats au supermarché peuvent se répartir entre les catégories \emph{Alimentation : Épicerie} et \emph{Soins : Habillement}). Vous pouvez aussi répartir une somme sur la même catégorie, mais avec des informations différentes dans le champ \menus{Remarques} (par exemple pour mémoriser le détail d'une opération).

\vspacepdf{3mm}
\Note{}: please ensure you follow the syntax correctly, leaving a space before AND after the \enquote{:} between the category and a new sub-category.
%\Note{}: veuillez à bien respectez la syntaxe avec un espace avant ET après les \frquote{:} entre la catégorie et une nouvelle sous-catégorie.


\subsection{Entering the breaken-down transaction\label{transactions-breakdown-master}}
%\subsection{Saisie de l'opération ventilée\label{transactions-breakdown-master}}

To \indexword{enter a breaken-down transaction}\index{transaction!breaken-down}, proceed as follows:
%Pour \indexword{saisir une opération ventilée}\index{opération!ventilée}, procédez comme suit:
	 
\begin{enumerate}
	\item in the form, start entering your transaction as you would for a normal transaction until you reach the \menus{Categories} field, where you should select the \menus{Split Transaction} category (towards the end of the drop-down menu, but it’s easier to type the letters \enquote{sp});
	%\item dans le formulaire, commencez la saisie de votre opération comme une opération normale jusqu'au champ \menus{Catégories}, où vous choisissez la catégorie \menus{Opération ventilée} (vers la fin du menu déroulant, mais il est plus simple de taper la lettre %\frquote{o});
	\item continue entering the remaining details, then validate:
	%\item continuez à saisir les autres données, puis validez:
	\item the following appear in the list of transactions:
	%\item dans la liste des opérations, apparaissent:
	\begin{itemize}
		\item the first line of the new operation is written in \textcolor{red}{\textsf{red}} and is therefore highlighted; it will remain in \textcolor{red}{\textsf{red}} as long as the total of the sub-transactions differs from the amount of the transaction;
		%\item la première ligne de la nouvelle opération écrite en \textcolor{red}{\textsf{rouge}} donc sélectionnée, et qui restera écrite en \textcolor{red}{\textsf{rouge}} tant que le total des sous-opérations sera différent du montant de l'opération;
		\item and below that, a new line, with \fboxsep=3pt{\colorbox[RGB]{245,155,155}{dark pink background}}, with a total of zero and a difference equal to the difference between this total and the total value of the split transaction, i.e. a negative figure\refimage{transactions-breakdown-img};
		%\item et en-dessous une nouvelle ligne, sur \fboxsep=3pt{\colorbox[RGB]{245,155,155}{fond rose foncé}}, avec un total égal à zéro et un écart égal à la différence entre ce total et la valeur totale de l'opération ventilée, donc en négatif\refimage{transactions-breakdown-img};
	\end{itemize}

		% image centrée 
		\begin{figure}[htbp]
		\begin{center}
		\includegraphics[width=0.95\textwidth]{image/screenshot/transactions_breakdown}
		\end{center}
		\vspace{-20pt}
		\caption{Entering the split transaction}
		%\caption{Saisie de l'opération ventilée}
		\label{transactions-breakdown-img}
		\vspace{-10pt}
		\end{figure}
		% image centrée

	\item in the entry form, the cursor is in the \menus{Debit} field, and the \menus{Date}, \menus{Payee} and \menus{Payment method} fields are greyed out, meaning they can no longer be edited: this is the sub-transaction form, which is ready for entering sub-transactions;
	%\item dans le formulaire de saisie, le curseur est dans le champ \menus{Débit}, et les champs \menus{Date}, \menus{Tiers} et \menus{Moyen de paiement} sont grisés, ce qui signifie qu'on ne peut plus les modifier: il s'agit du formulaire de sous-opération, qui est prêt pour la saisie des sous-opérations;
	\item go to the next subsection \menus{Entering split sub-transactions}.
	%\item passez à la sous-section suivante \menus{Saisie des sous-opérations ventilées}.
\end{enumerate}


\subsection{Entering split sub-transactions\label{transactions-breakdown-slave}}
%\subsection{Saisie des sous-opérations ventilées\label{transactions-breakdown-slave}}

\Note{}: the amount of a sub-transaction may be the result of a transaction involving several purchases or expenses listed on a receipt or invoice; you can therefore use the \indexword{calculator}\index{calculator} in the \menus{Debit} or \menus{Credit} input fields to more easily obtain the amounts of sub-transactions (see the section \vref{transactions-new-numbers}, \menus{Entering formulas}).
%\Note{}: le montant d'une sous-opération peut-être le résultat d'une opération portant sur plusieurs achats ou dépenses inscrites sur un ticket de caisse ou une facture; vous pouvez donc utiliser la \indexword{calculatrice}\index{calculatrice} dans les champs de saisie \menus{Débit} ou \menus{Crédit} pour obtenir plus facilement les montants des sous-opérations (voir la section \vref{transactions-new-numbers}, \menus{Saisie de formules}).

\vspacepdf{2mm}
To enter \indexword{split sub-transactions}\index{sub-transactions!split}, continue typing as follows:
%Pour saisir les \indexword{sous-opérations ventilées}\index{sous-opération!ventilée}, continuez la saisie comme suit:
%\vspacepdf{3mm}
\begin{enumerate}
	\item in this sub-operation form, fill in the empty fields (those not greyed out), then save your first sub-operation;
	%\item dans ce formulaire de sous-opération, saisissez les champs libres (non grisés), puis validez votre première sous-opération;

		% image centrée 
		\begin{figure}[htp]
		\begin{center}
		\includegraphics[width=0.95\textwidth]{image/screenshot/transactions_breakdown_sub1}
		\end{center}
		\vspace{-20pt}
		\caption{Entering the first split sub-transaction}
		%\caption{Saisie de la première sous-opération ventilée}
		\label{transactions-breakdown-sub1-img}
		\vspace{-10pt}
		\end{figure}
		% image centrée
 
	\item in the list of operations, this sub-transaction is displayed on \fboxsep=3pt{\colorbox[RGB]{240,200,200}{pink background}}, and a new line appears, on \fboxsep=3pt{\colorbox[RGB]{245,155,155}{dark pink background}} and is therefore selected, displaying a total equal to the first sub-transaction entered, in negative figures, and a difference equal to the difference between the total value of the split transaction and the value of this sub-transaction, in negative form\refimage{transactions-breakdown-sub1-img};
	%\item dans la liste des opérations, cette sous-opération s'affiche sur \fboxsep=3pt{\colorbox[RGB]{240,200,200}{fond rose}}, et une nouvelle ligne apparaît, sur \fboxsep=3pt{\colorbox[RGB]{245,155,155}{fond rose foncé}} donc sélectionnée, qui affiche un total égal à la première sous-opération saisie, en négatif, et un écart égal à la différence entre la valeur totale de l'opération ventilée et la valeur de cette sous-opération, en négatif\refimage{transactions-breakdown-sub1-img};
	\item there are at least two sub-transactions in a split transaction (otherwise there is no point in doing it\ldots); therefore, in the sub-transaction form, enter a second sub-transaction in the same way, then click on \menus{Validate}: this sub-transaction appears on \fboxsep=3pt{\colorbox[RGB]{240,200,200}{pink background}} and a new line appears, on \fboxsep=3pt{\colorbox[RGB]{245,155,155}{dark pink background}} and is therefore selected, displaying the latest total of the sub-transactions entered and the latest difference \refimage{transactions-breakdown-sub2-img};
	the total of the sub-transactions has increased (negatively) by the amount of this second sub-transaction, whilst the difference has decreased (negatively) by the same amount;
	%\item il y a au moins deux sous-opérations dans une opération ventilée (sinon cela n'a aucun intérêt\ldots); donc, dans le formulaire de sous-opération, saisissez une deuxième sous-opération de la même manière, puis validez: cette sous-opération s'affiche sur \fboxsep=3pt{\colorbox[RGB]{240,200,200}{fond rose}} et une nouvelle ligne apparaît, sur \fboxsep=3pt{\colorbox[RGB]{245,155,155}{fond rose foncé}} donc sélectionnée, qui affiche le dernier total des sous-opérations saisies et le dernier écart \refimage{transactions-breakdown-sub2-img}; le total des sous-opérations a augmenté (en négatif) du montant de cette deuxième sous-opération, tandis que l'écart a diminué (en négatif) du même montant; 

	% image centrée 
	\begin{figure}[htp]
	\begin{center}
	\includegraphics[width=0.95\textwidth]{image/screenshot/transactions_breakdown_sub2}
	\end{center}
	\vspace{-15pt}
	\caption{Entering the second split sub-transaction}
	%\caption{Saisie de la deuxième sous-opération ventilée}
	\label{transactions-breakdown-sub2-img}
	\vspace{-10pt}
	\end{figure}
	% image centrée 
\newpage
	\item if the difference is not zero, a third sub-transaction can then be entered in the same way\refimage{transactions-breakdown-sub3-img},
	%\item si l'écart n'est pas nul, on peut alors saisir une troisième sous-opération de la même manière\refimage{transactions-breakdown-sub3-img},
	
	% image centrée 
	\begin{figure}[h]
	\begin{center}
	\includegraphics[width=0.95\textwidth]{image/screenshot/transactions_breakdown_sub3}
	\end{center}
	\vspace{-10pt}
	\caption{Entering the third split sub-transaction} 
	%\caption{Saisie de la troisième sous-opération ventilée}
	\label{transactions-breakdown-sub3-img}
	\vspace{-5pt}
	\end{figure}
	% image centrée 

	\item and so on until the amount you have entered for your last sub-transaction in the form is exactly equal to the difference shown.
	%\item et ainsi de suite jusqu'à ce que le montant saisi de votre dernière sous-opération dans le formulaire soit exactement égal à l'écart affiché.
	
	\Note{}: until the transaction is complete, it remains displayed in \textcolor{red} red.
	%\Note{}: tant que l'opération n'est pas complète, celle-çi reste écrite en \textcolor{red}{rouge}.
\end{enumerate}


\subsection{Validation of the split transaction\label{transactions-breakdown-validation}}
%\subsection{Validation de l'opération ventilée\label{transactions-breakdown-validation}}

Once this final sub-operation has been validated, your split transaction will appear in the list of transactions just as a normal transaction would, and the active line, highlighted in \fboxsep=3pt{\colorbox[RGB]{245,155,155}{dark pink background}}, will move to the last empty line in the list.
%À la validation de cette dernière sous-opération, votre opération ventilée s'affiche alors dans la liste des opérations comme pour une opération normale, et la ligne active, sur \fboxsep=3pt{\colorbox[RGB]{245,155,155}{fond rose foncé}}, passe à la dernière ligne vide de la liste.

In the list of transactions, you will be able to identify your transaction by the presence of a triangle to the left of the first column \refimage{transactions-breakdown-display-img}.
%Dans la liste des opérations, vous pourrez reconnaitre votre opération ventilée par la présence d'un triangle à gauche de la première colonne \refimage{transactions-breakdown-display-img}.

\begin{figure}[!htbp]
	%\vspace{-15pt}
	\begin{center}
		\includegraphics[width=0.95\textwidth]{image/screenshot/transactions_breakdown_display}
	\end{center}
	\vspace{-10pt}
	\caption{Split transaction in the list of transactions}
	%\caption{Opération ventilée dans la liste des opérations}
	\label{transactions-breakdown-display-img}
	\vspace{-5pt}
\end{figure}


\subsection{Displaying the split transaction\label{transactions-breakdown-display}}
%\subsection{Affichage de l'opération ventilée\label{transactions-breakdown-display}}

By clicking on the triangle to the left of the split transaction, you select that transaction and view the details of its sub-transactions \refimage{transactions-breakdown-fulldisplay-img}.
%En cliquant sur le triangle à gauche de l'opération ventilée, vous sélectionnez cette opération et vous affichez le détail de ses sous-opérations \refimage{transactions-breakdown-fulldisplay-img}.

\begin{figure}[htbp]
	%\vspace{-15pt}
	\begin{center}
		\includegraphics[width=0.95\textwidth]{image/screenshot/transactions_breakdown_fulldisplay}
	\end{center}
	\vspace{-10pt}
	\caption{Full view of the split transaction}
	%\caption{Affichage complet de l'opération ventilée}
	\label{transactions-breakdown-fulldisplay-img}
	\vspace{-5pt}
\end{figure}

\vspacepdf{3mm}
\Note{}: depending on the theme of your desktop environment or the window manager you are using, these triangles can be replaced with other characters such as +, -, >, <, etc.
%\Note{}: ce triangle peut être remplacé, en fonction du thème de l'environnement de bureau ou du gestionnaire de fenêtres que vous utilisez, par d'autres caractères tels que +, -, >, <, etc.


\subsection{Automatic completion of the form for a split transaction\label{transactions-breakdown-autofill}}%\label{transactions-breakdown-recall}
%\subsection{Remplissage automatique du formulaire d'une opération ventilée\label{transactions-breakdown-autofill}}%\label{transactions-breakdown-recall}

The \indexword{automatic completion}\index{split transaction!autofill} feature for a split transaction works in the same way as for a standard transaction (see the section \vref{transactions-new-autofill}, \menus{Automatic form completion}).
%Le \indexword{remplissage automatique}\index{opération ventilée!remplissage automatique} du formulaire d'une opération ventilée fonctionne comme pour celui d'une opération normale (voir la section \vref{transactions-new-autofill}, \menus{Remplissage automatique du formulaire}).

However, for a broken-down operation, a tick box labelled \menus{Recover the children} appears at the bottom of the input form, to the left of the \menus{Cancel} button, and you can tick this box as required.
%Cependant, pour une opération ventilée, une case à cocher intitulée \menus{Restaurer les sous-opérations} apparaît en bas du formulaire de saisie, à gauche du bouton \menus{Annuler}, et vous pouvez cocher cette case selon votre besoin.


\section{Editing an transaction\label{transactions-modify}}
%\section{Modification d'une opération\label{transactions-modify}}

To edit an transaction, a broken-down transaction, or one of its sub-transactions, select it from the list of transactions, where it will be displayed in \fboxsep=3pt{\colorbox[RGB]{245,155,155}{dark pink background}}; it will also appear in the entry form, where you can select the desired fields by \texttt{left-click} and modify the values as you see fit.
%Pour modifier une opération, ou une opération ventilée ou une de ses sous-opérations, sélectionnez-la dans la liste des opérations où elle s'affichera sur \fboxsep=3pt{\colorbox[RGB]{245,155,155}{fond rose foncé}}; elle s'affichera également dans le formulaire de saisie, où vous pourrez sélectionner les champs désirés par un \texttt{clic-gauche} et modifier les valeurs à votre convenance.

\vspace{2mm}
If the entry form is hidden, do one of the following to display it:
%Si le formulaire de saisie est masqué, procédez d'une des manières suivantes pour l'afficher:
\begin{itemize}
	\item in the menu bar, select \menus{Edit - New transaction} if an account is selected;
	%\item dans la barre de menu, sélectionnez \menus{Éditer - Nouvelle opération} si un compte est sélectionné;
	\item in the menu bar, select \menus{View} and tick the box next to the \menus{Show transaction form} option;
	%\item dans la barre de menu, sélectionnez \menus{Affichage} et cochez la case de l'option \menus{Montrer le formulaire de saisie des opérations};
	\item select the blank line at the bottom of the list of transactions and press the \keys{Enter \return} key;
	%\item sélectionnez la ligne vide en bas de la liste des opérations et appuyez sur la touche \keys{Entrée \return};
	\item double-click on the empty line at the bottom of the list of transactions;
	%\item double-cliquez sur la ligne vide en bas de la liste des opérations;
	\item press the \keys{Ctrl+T} key combination;
	%\item appuyez sur la combinaison de touches \keys{Ctrl+T};
	\item in the toolbar, click on \menus{New transaction};
	%\item dans la barre d'outils, cliquez sur \menus{Nouvelle opération};
	\item in the context menu of the transaction list, select \menus{New transaction};
	%\item dans le menu contextuel de la liste des opérations, sélectionnez \menus{Nouvelle opération};
	\item click on the small triangle next to the label \indexword{show/hide} on the transaction entry form to the left of the label.
	%\item cliquez sur le petit triangle d'\indexword{affichage/masquage}\index{formulaire de saisie!masquer-afficher} du formulaire de saisie des opérations à gauche de son libellé.
	
	% saut de ligne pour indentation correcte de la note dans la liste
	\vspacepdf{2mm}
	\Note{}: depending on the theme of your desktop environment or the window manager you are using, these triangles can be replaced with other characters such as +, -, >, <, etc.
	%\Note{}: ce triangle peut être remplacé, en fonction du thème de l'environnement de bureau ou du gestionnaire de fenêtres que vous utilisez, par d'autres caractères tels que +, -, >, <, etc.
\end{itemize}
  
When the process is complete, the date appears in \fboxsep=3pt{\colorbox[RGB]{61,174,233}{\textcolor{white}{\textsf{white on a blue background}}}} in the form.
%A l'issue, la date s'affiche en \fboxsep=3pt{\colorbox[RGB]{61,174,233}{\textcolor{white}{\textsf{blanc sur fond bleu}}}} dans le formulaire.

You can then edit any information you wish by navigating through the entry fields.
%Il est alors possible de modifier toutes les informations désirées, en se déplaçant dans les champs de saisie.

The \keys{Ctrl+P} key combination sets or clears the \indexword{pointed transaction}\index{transaction!pointed} status for a selected transaction, and \keys{Ctrl+R} sets or clears the \indexword{reconciled transaction}\index{transaction!reconciled} status (see the \vref{reconciliation} chapter, \menus{Bank Reconciliation}). You can also assign a transaction the status of a pointed transaction—but only this status—by using the \keys{Ctrl} + \texttt{Left Click} shortcut in the \enquote{P/R} column on the line of the transaction to be pointed.
%La combinaison de touches \keys{Ctrl+P} donne ou retire à une opération sélectionnée le statut d'\indexword{opération pointée}\index{opération!pointée}, et \keys{Ctrl+R} celui d'\indexword{opération rapprochée}\index{opération!rapprochée} (voir le chapitre \vref{reconciliation}, \menus{Rapprochement bancaire}). Vous pouvez aussi donner à une opération le statut d'opération pointée, mais uniquement celui-ci, par la combinaison \keys{Ctrl+Clic} dans la colonne \frquote{P/R}, sur la ligne de l'opération à pointer.

% espace pour changement de thème
\vspacepdf{3mm}%TODO: En ATTENTE
If the transaction has already been reconciled (see the section \vref{reconciliation}, \menus{Bank reconciliation}), you cannot change its amount. The other fields remain editable. If you absolutely must change the amount (though if you have already reconciled the transaction, this should not be necessary), proceed as follows:
%Si l'opération a déjà été rapprochée (voir le chapitre \vref{reconciliation}, \menus{Rapprochement bancaire}), il est par contre impossible de modifier son montant. Les autres champs restent modifiables. Si vous avez vraiment et impérativement besoin de modifier le montant (mais si vous avez déjà rapproché l'opération, cela ne devrait pas être le cas), procédez comme suit:

\begin{enumerate}
	\item select the transaction;
	%\item sélectionnez l'opération;
	\item cancel its \emph{Reconciled} status by pressing the \keys{Ctrl+R} key combination;
	%\item annulez son statut \emph{Rapproché} en appuyant sur la combinaison de touches \keys{Ctrl+R};
	\item change the amount;
	%\item modifiez le montant;%TODO verify
	\item press the \keys{Ctrl+R} key combination again. If you do not do this, the reconciliation for this transaction will be inaccurate by an amount equal to the difference between the previous and new amounts for this transaction.
	%\item appuyez à nouveau sur la combinaison de touches \keys{Ctrl+R}. Si vous ne le faisiez pas, le rapprochement relatif à cette opération serait faussé, d'une valeur égale à la différence entre le précédent et le nouveau montant de cette opération.
\end{enumerate}

\Attention{}: cette manipulation est hautement déconseillée, car elle fausse momentanément le rapprochement associé à cette opération, et il redevient exact à la fin, si vous n'en avez pas modifié le montant (cette manipulation sera d'ailleurs probablement automatiquement désactivée dans une future version de Grisbi dans le cas d'une comptabilité de type association).


\section{Deleting a transaction\label{transactions-delete}}
%\section{Suppression d'une opération\label{transactions-delete}}

To \indexword{delete a transaction}\index{transaction!delete}, select it, and use one of the following methods:
%Pour supprimer une opération, sélectionnez-la, et utilisez l'une de ces méthodes:
\begin{itemize}
	\item press the \keys{Delete} key on the keyboard;
	%\item appuyez sur la touche du clavier \keys{Suppr};
	\item in the menu bar, select \menus{Edit - Remove transaction};
	%\item dans la barre de menu, sélectionnez \menus{Éditer - Supprimer une opération};
	\item in the toolbar, \texttt{left-click} on the \menus{Delete} button;
	%\item dans la barre d'outils, cliquez-gauche sur \menus{Supprimer};
	\item after \texttt{right-click}, select \menus{Delete transaction} from the context menu.
	%\item après un \texttt{clic-droit}, dans le menu contextuel, sélectionnez \menus{Supprimer l'opération}.
\end{itemize}

% espace pour changement de thème
\vspacepdf{2mm}
Deleting a broken-down sub-transaction is done in the same way, and the broken-down transaction will turn \textcolor{red}{red}. You must therefore ensure that the sum of the sub-transactions ultimately equals the value of the broken-down transaction.
%La suppression d'une sous-opération ventilée se fait de la même façon et l'opération ventilée passera en \textcolor{red}{rouge}. Vous devrez donc faire en sorte que finalement la somme des sous-opérations soit égale à la valeur de l'opération ventilée.

% espace pour changement de thème
\vspacepdf{2mm}
Deletion is not possible if:
%La suppression n'est pas possible si:
\begin{itemize}
	\item the transaction is open for editing: you must first exit edit mode by pressing the \keys{Esc} key;
	%\item l'opération est ouverte en modification: il faut préalablement sortir du mode édition, en utilisant la touche \keys{Échap};
	\item the transaction is \indexword{reconciled}\index{transaction!reconciled} (you must first deselect its \emph{Reconciled} status by pressing the \keys{Ctrl+R} key combination).
	%\item l'opération est \indexword{rapprochée}\index{opération!rapprochée} (il faut préalablement annuler son statut \emph{Rapproché} en appuyant sur la combinaison de touches \keys{Ctrl+R}).
\end{itemize}

% espace avant Attention ou Note: 5 mm
\vspacepdf{3mm}
\Attention{}: this final step is strongly discouraged as it distorts the reconciliation associated with this transaction.
%\Attention{}: cette dernière manipulation est hautement déconseillée car elle fausse le rapprochement associé à cette opération.


\section{Transaction selected as a template\label{transactions-model}}
%\section{Opération sélectionnée comme modèle\label{transactions-model}}

To use an \indexword{transaction as a template}\index{transaction!template}, select it, and use one of the following methods:
%Pour utiliser une \indexword{opération comme modèle}\index{opération!modèle}, sélectionnez-la, et utilisez l'une de ces méthodes:
\begin{itemize}
	\item in the menu bar, select \menus{Edit - Use selected transaction as a template};
	%\item dans la barre de menu, sélectionnez \menus{Éditer - Utiliser l'opération sélectionnée comme modèle};
	\item after \texttt{right-clicking}, select \menus{Use selected transaction as a template} from the context menu.
	%\item après un \texttt{clic-droit} sur une opération, dans le menu contextuel, sélectionnez \menus{Utiliser l'opération sélectionnée comme modèle}.
\end{itemize}

% espace pour changement de thème
\vspacepdf{2mm}
A new transaction with today's date is displayed in \fboxsep=3pt{\colorbox[RGB]{245,155,155}{dark pink background}}, on a new line below the last transaction; naturally, it has neither the \emph{Pointed} status nor the \emph{Reconciled} status.
%Une nouvelle opération avec la date du jour s'affiche sur \fboxsep=3pt{\colorbox[RGB]{245,155,155}{fond rose foncé}}, dans une nouvelle ligne en-dessous de la dernière opération; elle n'a évidemment ni le statut \emph{Pointé} ni le statut \emph{Rapproché}.

You can then edit this new transaction, for example by changing its date.
%Vous pouvez alors éditer cette nouvelle opération, par exemple en modifiant sa date.


\section{Cloning a transaction\label{transactions-duplicate}}
%\section{Clonage d'une opération\label{transactions-duplicate}}

To \indexword{clone a transaction}\index{transaction!clone}, select it, then use one of the following two methods:
%Pour cloner une opération, sélectionnez-la, puis utilisez l'une de ces deux méthodes:
\begin{itemize}
	\item in the menu bar, select \menus{Edit - Clone transaction};
	%\item dans la barre de menu, sélectionnez \menus{Éditer - Cloner l'opération};
	\item after a \texttt{right-click}, select \menus{Clone transaction} from the context menu.
	%\item après un \texttt{clic-droit} sur une opération, dans le menu contextuel, sélectionnez \menus{Cloner l'opération}.
\end{itemize}

% espace pour changement de thème
\vspacepdf{2mm}
A new transaction with the previously selected transaction date appears in \fboxsep=3pt{\colorbox[RGB]{245,155,155}{dark pink background}}, and is displayed on a new line below the selected transaction; naturally, it has neither the \emph{Pointed} status nor the \emph{Reconciled} status.
%Une nouvelle opération avec la date de l'opération sélectionnée précédemment apparaît sur \fboxsep=3pt{\colorbox[RGB]{245,155,155}{fond rose foncé}}, et s'affiche dans une nouvelle ligne en-dessous de l'opération sélectionnée; elle n'a évidemment ni le statut \emph{Pointé} ni le statut \emph{Rapproché}.

You can then edit this new transaction, for example by changing its date.
%Vous pouvez alors éditer cette nouvelle opération, par exemple en modifiant la date.


\section{Converting a transaction into a scheduled transaction\label{transactions-schedule}}
%\section{Conversion d'une opération en opération planifiée\label{transactions-schedule}}

To \indexword{convert an transaction into a scheduled transaction}\index{transaction!into a scheduled transaction}, select it, then use one of these two methods:
%Pour \indexword{convertir une opération en opération planifiée}\index{opération!en opération planifiée}, sélectionnez-la, puis utilisez l'une de ces deux méthodes:
\begin{itemize}
	\item in the menu bar, select \menus{Edit - Convert transaction to scheduled transaction};
	%\item dans la barre de menu, sélectionnez \menus{Éditer - Convertir en opération planifiée};
	\item after \texttt{right-clicking}, select \menus{Convert transaction to scheduled transaction} from the context menu.
	%\item après un \texttt{clic-droit}, dans le menu contextuel, sélectionnez \menus{Convertir l'opération en opération planifiée}.
\end{itemize}

% espace pour changement de thème
\vspacepdf{2mm}
The \menus{Scheduler} tab opens. The new transaction appears in the list of scheduled transactions, and the entry form is displayed with the transaction settings. Check or edit these settings as required, then confirm.
%L'onglet \menus{Échéancier} s'affiche. La nouvelle opération apparaît dans la liste des opérations planifiées, et le formulaire de saisie s'affiche avec les paramètres de l'opération. Vérifiez ou éditez ces paramètres selon vos besoins, puis validez.


\section{Moving a transaction to another account\label{transactions-move}}
%\section{Déplacement d'une opération vers un autre compte\label{transactions-move}}

To \indexword{move a transaction to another account}\indexword{transaction!move to another account}, select it, then use one of these two methods:
%Pour \indexword{déplacer une opération vers un autre compte}\indexword{opération!déplacer vers un autre compte}, sélectionnez-la, puis utilisez l'une de ces deux méthodes:
\begin{itemize}
	\item in the menu bar, select \menus{Edit - Move transaction to another account} and choose the destination account from the list;
	%\item dans la barre de menu, sélectionnez \menus{Éditer - Déplacer l'opération vers un autre compte} et choisissez le compte destinataire dans la liste;
	\item after \texttt{right-clicking}, select \menus{Move transaction to another account} from the context menu and choose the destination account from the list.
	%\item après un \texttt{clic-droit}, dans le menu contextuel, sélectionnez \menus{Déplacer l'opération vers un autre compte} et choisissez le compte destinataire dans la liste.
\end{itemize}

It goes without saying that both accounts must be in the same Grisbi created by file.
%Il est bien entendu que les deux comptes doivent se trouver dans le même fichier créé par Grisbi.

% espace avant Attention ou Note: 5 mm
\vspacepdf{3mm}
\Note{}: It is not possible to move a transaction if it has the status \emph{Reconciled}. However, it is possible if it has the status \emph{Pointed}.
%\Note{}: il est impossible de déplacer une opération si elle a le statut \emph{Rapproché}. Par contre, c'est possible si elle a le statut \emph{Pointé}.


\section{Nouveau tiers, catégorie ou imputation budgétaire\label{transactions-fillcombo}}

Il peut arriver que vous ayez besoin d'un nouveau tiers, d'une nouvelle catégorie ou d'une nouvelle imputation budgétaire pendant une saisie d'opération. Dans ce cas, Grisbi permet d'en faire la création %\frquote{à la volée}
: saisissez le nouveau nom dans la zone de liste déroulante, et Grisbi l'ajoutera automatiquement à la liste. Vous pourrez ultérieurement faire toutes les modifications nécessaires dans les onglets \menus{Tiers} (voir la section \vref{thirdparties-modify}, \menus{Modification d'un tiers}), \menus{Catégories} (voir la section \vref{categories-modify}, \menus{Modification d'une catégorie ou d'une sous-catégorie}) ou \menus{Imputations budgétaires} (voir la section \vref{budgetarylines-modify}, \menus{Modification d'une (sous-)imputation budgétaire}).
 
% espace avant Attention ou Note: 5 mm
\vspacepdf{3mm}
\Note{}: pour saisir une \indexword{catégorie}\index{catégories!syntaxe}\index{sous-catégories!syntaxe}\index{syntaxe!catégories : sous-catégories}
 (ou une \indexword{imputation budgétaire} : sous-imputation budgétaire\index{imputations budgétaires!syntaxe}\index{sous-imputations budgétaires!syntaxe}\index{syntaxe!imputations budgétaires : sous-imputations budgétaires}) avec une sous-catégorie (ou sous-imputation), il faut respecter la syntaxe \emph{catégorie : sous-catégorie} (avec les deux points comme séparateur, encadrés par des espaces), sinon Grisbi ne comprendra pas. De même si vous ajoutez une nouvelle sous-catégorie (ou sous-imputation budgétaire) à une catégorie (ou imputation) existante.
% espace après Attention ou Note: 5 mm
 \vspacepdf{3mm}

Vous pouvez également créer un tiers, une catégorie ou une imputation budgétaire directement dans leur onglet respectif, sélectionné dans le panneau de navigation.


\section{Impression de la liste des opérations\label{transactions-print}}
%TODO verify
Pour imprimer la liste des opérations, cliquez sur le bouton \menus{Imprimer} de la barre d'outils; une fenêtre d'impression s'ouvre, dont l'aspect et les fonctions dépendent de votre gestionnaire d'impression; vous aurez le plus souvent les choix suivants:

\begin{itemize}
	\item imprimer dans un fichier (au format \gls{PDF}, \gls{PostScript}, ou \gls{SVG});
	\item imprimer avec votre imprimante.
\end{itemize}

En fonction de votre gestionnaire d'impression, vous pourrez disposer de réglages divers tels que la taille et l'orientation de la feuille, la résolution, la police d'impression et sa taille, etc.

% espace avant Attention ou Note: 5 mm
\vspacepdf{3mm}
\Note{}: la liste des opérations de certains comptes peut être très longue; affichez un aperçu avant impression pour vérifier ce que vous allez imprimer.


%XXXXXXXXXXXXXXXXXXXXXXXXXXXXXXXXXXXXXXXXXXXXXXXXXXXXXXXXXXXXXXXXXXXXXXXXXXXXXXXXXX
%
% [benj] supprimé car franchement, ça n'a rien à faire ici.
%
% PLEASE NOTE! DO NOT TRANSLATE THE FOLLOWING SECTION
%\section{Quelques mots sur un anglicisme trop répandu}
%
%On voit régulièrement dans les documentations, livres et autres articles,
%publiés sur l'Internet et ailleurs, le mot \emph{complétion}. Ceci est un
%horrible anglicisme.
%
%Que l'on assimile un mot anglais lorsque le mot français équivalent n'existe pas
%peut se comprendre, encore que le néologisme intelligent me semble préférable.
%Il ne s'agit bien sûr pas de tomber dans le ridicule de nos académiciens qui
%traduisent mail par mèl et CD-Rom par cédérom. Passons sur leur manque
%d'imagination alors qu'ils ont su par autrefois créer des mots comme disquette,
%informatique et ordinateur, qui n'existent nulle part ailleurs!
%
%Dans le cas précis du mot \emph{complétion} un traducteur du \emph{Guide de
%l'administration réseau sous linux,} Editions O'Reilly, avait cru bien faire en
%essayant de créer le mot \emph{complétement} (avec un accent aigu) qui voulait
%signifier \emph{action de compléter}. L'initiative était louable mais ô combien
%risquée \ldots En effet lorsque l'on s'essaye à ce genre d'exercice de style, il est
%prudent de maîtriser tous les rouages de la langue. Et en l'occurrence il existe
%une règle dite la \emph{règle de l'ernance} (page 56 du \emph{Bescherelle
%Orthographe}) \footnote{Si quelqu'un possède ce livre je serais très heureux
%qu'il me fasse parvenir le texte exact de cette règle. Je ne l'ai 
%malheureusement plus\ldots} qui explique parfaitement le mécanisme de
%changement d'accent entre un verbe et le substantif dérivé
%\footm{André Cianfarini, qui se définit lui-même comme \emph{un développeur autodidacte qui n'a jamais fini sa thèse de grammaire pour cause de relations passionnelles avec l'informatique} me communique les précisions suivantes: 
%\begin{quote}
%C'est un processus phonologique qui est à l'origine de ces
%transformations; pour simplifier, on pourrait le résumer ainsi:
%un "e" fermé (l'avant-dernier "e" de "préférer" par exemple) devient un
%"e" ouvert (l'avant-dernier "e" de "je préfère" par exemple) lorsque la
%syllabe qui le suit est atone.
%\end{quote} 
%}. Si ce malheureux
%traducteur l'avait connue il n'aurait pas commis cet affreux barbarisme\ldots
%
%Et si avant toute chose ce traducteur avait eu l'idée de vérifier dans son
%dictionnaire que le mot n'existait pas déjà\ldots il aurait trouvé:
%
%\begin{description}
%
%\item [complètement]n.m. Action de compléter: \emph{le complètement
%d'un livre}.
%
%\end{description}
%
%Et j'ai bien pris soin de chercher ce mot dans le plus vieil exemplaire du
%\emph{Larousse universel} en ma possession, soit l'édition de 1949!
%
%Donc adoptons sans aucun remord le mot complètement, et, puisque cette action
%est automatique, allons sans crainte jusqu'à écrire auto-complètement, ce qui
%est parfaitement compréhensible, intuitif et 100\% français!
%
%Cocorico!
%
%Et puissent ces quelques lignes tordre définitivement le cou au mot complétion,
%cette \emph{affreuseté non-permissable}!;-)
%
%from now you can translate again